% ============================================================
%  H. Høffding, Relation som Kategori:
%  en erkendelsesteoretisk Undersøgelse
%  Det Kgl. Danske Videnskabernes Selskab,
%  Filosofiske Meddelelser I, 3.
%  København: Bianco Lunos Bogtrykkeri, 1921.
%
%  TRANSLATION (English)
%  Translator: [HH]
%  Status: IN PROGRESS. Source of truth is transcription.tex in this
%  directory (COMPLETE, printed pp. 3–109, made from the Royal Danish
%  Academy's digitisation of the 1921 printing).
%
%  The δ section (§3d.δ, pp. 72–77) was REDONE 2026-08-23 from that
%  transcription. Its previous version had been made from a pre-scan Danish
%  text that collation showed to diverge from the printing in eleven places;
%  the divergences that reach the English are itemised at the head of the
%  section itself.
%
% ------------------------------------------------------------
%  CONVENTIONS (follow these for every remaining section)
% ------------------------------------------------------------
%  * MARKERS in this file read "% [tr'anslation to be added]" (without the
%    apostrophe) — NOT the "% [te'xt to be added: pp. X--Y]" form the rest of
%    the repo uses. A grep for the repo-standard string reports this file as
%    clean when it is not; that is very likely how this book came to be
%    recorded as further along than it was. The apostrophes above are
%    deliberate, so that this header does not itself inflate the marker count.
%    To count what is genuinely left:
%      grep -c 'tr'"'"'anslation to be added' translation.tex   # header-safe
%    or just read the count off check.py / the batch report.
%  * Quotes. The Danish sets GUILLEMETS »…« throughout (this book does not use
%    the low-high „…“ pair). Render them ``…'' in English. Foreign work titles
%    that the printing puts in guillemets stay in quotes, NOT italic.
%  * SMALL CAPS are the footnote house style for proper names at first
%    mention, and the printing uses them heavily; carry every \textsc{} over
%    from the transcription. The companion volume's translation (Totalitet som
%    Kategori) has 353 of them — that is the calibration.
%  * Italic. The transcription distinguishes \emph{} (author's stress, only 8
%    sites in the whole book) from \textit{} (work titles and foreign phrases
%    set in italic by convention, 74 sites). Keep the distinction; do NOT
%    invent italic for foreign phrases the printing sets in roman — that was
%    the single largest error in the superseded δ.
%  * Letterspacing (spærret) is structural ONLY in this book — subsection
%    heads and the Indhold. It is never authorial stress in body text. So a
%    letterspaced run is a head, not an emphasis.
%  * NO page markers. Neither this file nor the companion volume's translation
%    carries "% --- p. N ---" comments; page ranges live in the section
%    comments only.
%  * Printer's defects that are Danish-glyph-level (e.g. "Belydning" for
%    "Betydning", p. 75 n. 1) cannot be carried into English. Translate the
%    intended sense and note the defect in a % comment at the site.
% ============================================================

\documentclass[12pt]{article}
\usepackage{fontspec}
\usepackage{unicode-math}
\setmainfont{Libertinus Serif}
\setmathfont{Libertinus Math}
\usepackage[english]{babel}
\usepackage[protrusion=true]{microtype}
\usepackage[margin=1.4in]{geometry}
\usepackage{setspace}
\usepackage{fancyhdr}
\usepackage[colorlinks=true,linkcolor=black,urlcolor=black,
  pdftitle={Relation as Category},
  pdfauthor={Harald Høffding}]{hyperref}
\setlength{\headheight}{15pt}
\pagestyle{fancy}
\fancyhf{}
\fancyhead[LE,RO]{\thepage}
\fancyhead[RE]{\textit{Relation as Category}}
\fancyhead[LO]{\textit{\rightmark}}
\setlength{\parindent}{1.5em}
\setlength{\parskip}{0pt}
\setstretch{1.3}

\renewcommand{\thesection}{\arabic{section}.}
\renewcommand{\thesubsection}{\alph{subsection}.}

\title{\textbf{Relation as Category}\\[0.5em]
  \large An Epistemological Investigation}
\author{Harald Høffding}
\date{Translated from Danish by [HH]\\[0.4em]
  \textit{Det Kgl.\ Danske Videnskabernes Selskab,
  Filosofiske Meddelelser} I, 3.\\
  Copenhagen: Bianco Lunos Bogtrykkeri, 1921}

\begin{document}
\maketitle
\tableofcontents
\bigskip


% ============================================================
%  1. THE HISTORY OF THE CONCEPT OF RELATION
% ============================================================

\section{The History of the Concept of Relation}
\label{sec:1}
% pp. 3--16
% Range corrected from the skeleton's "3--15": section 1's closing paragraph
% ("That all necessity ... the work of thought.") runs on across the page turn
% and the section 2 head does not open until roughly two fifths down p. 16.

\noindent The present treatise is, like an earlier treatise on ``Totalitet
som Kategori'' (Videnskabernes Selskabs Skrifter, 1917), an attempt at the
further working out of a single point in my book ``Den menneskelige Tanke''
(1910), in which it was my endeavour to give a presentation of the principal
forms under which human thinking works, and, in connexion with these, of the
principal problems on which it works. The connexion between forms of thought
and problems of thought shows itself plainly in this, that the forms of
thought lie at the basis of the questions that are put --- questions which,
when they are put in a definite and sharp form, become problems.

After interest in epistemology had been awakened anew at the close of the
last century, it was above all the first fundamental propositions of human
knowledge that drew attention to themselves. Their possible grounding or
derivation was discussed, and so was the possibility of carrying them
completely through in every domain. But such fundamental propositions
contain certain fundamental concepts, which are applied more or less
involuntarily, partly by so-called sound common sense (sens commun), the
pre-scientific thinking, partly by scientific thinking, as the latter
displays itself in the history of the sciences. Both in the questions
thought puts, and in the answers it acknowledges, certain forms display
themselves, which it is the task of the doctrine of the categories, of
descriptive epistemology, to bring out and to order. The two principal
sources for the doctrine of the categories are the psychological study of
``sound common sense'' and the historical study of the development of
science proper through its various stages; and these sources stand in
constant reciprocal action with one another, in that science develops
successively out of the involuntary reflection of ordinary consciousness,
especially when a doubt has arisen, and in that, on the other hand, the
train of thought and the results of science slowly gain influence on the
form and content of ordinary consciousness, so that sound common sense is
not the same at all times.

While the earlier treatise dealt with the concept of totality and the
tendencies of thought connected with it, the present treatise shall seek to
throw light on the concept of relation and the principles of inquiry that
hang together with it. The two concepts point reciprocally to one another.
The concept of totality asserts itself along two different paths. All
knowledge consists in a comprehending together, an involuntary striving to
form wholes where such are not immediately given. But the immediately given
itself, the objects before us, can appear (and on closer investigation
perhaps always does appear) as wholes, and they can be understood only by
analysis of the connexion between their elements and with other wholes and
elements. Human knowledge forms wholes where such are not immediately given,
and dissolves them where they immediately lie before us. In both cases the
concept of the whole shows itself to hang together with the concept of
relation. Wholes can be formed only by objects being brought into definite
relations to one another, and given wholes are understood only by the
demonstration of the inner and outer relations that are decisive for their
subsistence. Relation therefore shows itself to be an essential form of
thought in every domain. To think is to find definite relations, to set one
object in relation to another; and the most general relations show
themselves to be continuity and discontinuity, likeness and difference.
These general or fundamental categories assume more special forms in the
formal categories (identity, analogy, negation, rationality), in the real
categories (causality, totality, development) and in the ideal categories
(the concepts of value). The closer demonstration of this is given in ``Den
menneskelige Tanke''.

The concept of relation will now in what follows receive illumination
through an investigation of the way in which it asserts itself in the
setting up and the discussion of the principal philosophical problems. In my
University Programme of 1902 and in ``Den menneskelige Tanke''
(p.~246---254) I have maintained that there are four such problems: the
psychological, the epistemological, the ethical and the cosmological, and
have pointed out the analogies between them. ---

\bigskip
An overview of the history of the doctrine of the categories (see ``Den
menneskelige Tanke'' p.~140---153) shows a decided tendency to place the
concept of relation in the first rank. In Aristotle, the first who set up a
doctrine of the categories, relation (πρὸς τι) is, to be sure, only one
among ten categories; but closer consideration shows that most (and at
bottom all) the other Aristotelian categories express particular species of
relation. Aristotle himself names as examples of relation: equality and
inequality, cause and effect, knowledge and the known. But in later Greek
thinking relation came to play an altogether decisive part, and here the
overview given in ``Den menneskelige Tanke'' needs an important
supplementing. Characteristic of this later Greek thinking, before mysticism
made its influence felt, was the critical-sceptical direction it struck out
upon. Plato stands, in spite of his speculative striving, as a precursor of
this direction, partly by the objections he allows to be voiced against the
doctrine of Ideas in ``Parmenides'', partly in the weight he lays, in later
dialogues, on knowledge expressing itself in the judgment, that is, in the
determination of the mutual relations of concepts (cf.\ my ``Bemærkninger om
den platoniske Dialog Parmenides'' in Videnskabernes Selskabs filosofiske
Meddelelser, 1920). But while Plato was, in spite of everything, convinced
that the truth had been found and was expressed in a world of eternal Ideas,
the later Academic school came to the result that it had not merely not been
found, but could not be found; what one could attain by the rational use of
experience was only probability. The most considerable man of this school is
Carneades (2nd cent.\ B.C.), who has rightly been called the mightiest mind
in Greece in these centuries. It became of decisive significance for the
concept of relation that he raised the question wherein our criterion for
truth and reality consists. As against the Stoics, who rested content with
the view that there were presentations in which the truth could be grasped
quite immediately (φαντασίαι καταληπτικαί)\footnote{On the interpretation of
this expression see \textsc{Paul Barth}: \textit{Die Stoa}, p.~66 f.}, he
maintains that, since several mutually conflicting presentations can appear
with this stamp, one cannot rely on any single presentation whatever. The
validity of a presentation can rest only on its agreement with other
presentations. It must not merely be ἀπερίστατος, not contradicted by other
presentations, but it must find confirmation through its connexion with
other presentations, just as physicians in their diagnoses do not hold to
any single mark, but rely on the connexion between all the marks (συνδρομὴ
τῶν σημείων). Further, there must in all important cases be required an
``analysing review'' of all the relations that have to do with presentations
(so that the presentation becomes a φαντασία
διεξωδευμένη).\footnote{\textit{Sextus Empiricus} ed.\ Bekker,
p.~225---232.}

It would be wrong to call Carneades a sceptic. He does not believe that any
presentation whatever has a claim to validity apart from its relation to
other presentations; but he thereby really only points to the path which
so-called sound common sense involuntarily takes when a doubt arises, and he
points to the method of empirical science. In his account of Carneades'
doctrine, Sextus compares his method with that of the empirical physicians,
and the conjecture has also been expressed that Carneades stands in debt to
this school of medicine.\footnote{\textsc{Victor Brochard}: \textit{Les
Sceptiques Grecs}. 1889, p.~183.} In the criterion of truth maintained by
Carneades the concept of relation plainly displays itself, and moreover in a
way already hinted at by Plato (Republic X, p.~602 C---603 A), as we shall
have occasion to mention at a later place.

If now the criterion of truth is the contradiction-free relation between
well-tested and thoroughly thought-out presentations, then a task is thereby
set that can never be absolutely solved. There are no limits to the number
of relations in which a presentation must be investigated in order that it
may stand its test. It shows itself here, what will show itself again and
again in what follows, that the concept of relation conditions at one and
the same time a limit and a task. By reason of this concept's predominant
significance in our knowledge, every understanding that is won expresses
certain definite relations and holds only for them, is to that extent
``relative''. But relation as category leads, on the other hand, always
beyond every understanding won, in that it points to the importance or the
necessity of seeking new relations which can make the determination of the
object understood more complete, and thereby the understanding itself more
deep-going. What offers resistance and arrests or limits thought, precisely
thereby sets it new tasks, yields new impulses. Charles Renouvier has
therefore rightly spoken of relation as method (la méthode des relations).
In the history of the concept of relation, right down to the most recent
time, the two tendencies contend --- limit and task, resistance and work,
conclusion and extension. And both tendencies spring from relation as
category.

Already in antiquity this contention shows itself. The Academics after
Carneades' time were mostly taken up with the conclusion, and thereby
approached the Stoics and their belief in a truth given once for all. The
Sceptics, on the other hand, held that when one can never have done with
winning the full truth, one has no right to assert that it is. Besides, it
may of course be asked by what criterion we can decide that the criterion
maintained by Carneades is the right one. --- While the Dogmatists reject
the concept of relation and cling to certain presentations as definitive,
the Sceptics maintain, on the basis of the concept of relation, that there
is no truth, for a relative truth is no truth.

The most eminent of these Sceptics, Aenesidemus (probably in the first
cent.\ B.C.), has given a systematic account of the various relations which
in his opinion make true knowledge impossible. He points to the differences
of living beings, particularly of men, as regards the faculty of
apprehension, --- to the difference between our various senses among
themselves, whereby ``one and the same thing'' comes to stand as different
for different senses, --- to the way in which our states and our standpoints
in space and time, particularly the distances from the object, condition our
apprehension, --- and to the influence of custom and law. In all he sets up
10 species of relation which render our knowledge uncertain. Sextus
Empiricus (ed.\ Bekker, p.~10) shows that these 10 species (τρόποι) can be
led back to three, in that either the apprehending and judging subject (ὁ
κρίνων), or the apprehended and judged object (τὸ κρινόμενον), or else both
at once, can determine the apprehension. And these three species of relation
he leads back again to the general concept of relation (ὁ πρὸς τι τρόπος).
Here for the first time relation is set up as a general category.

It has been wrongly supposed that in the part which the discussion of the
criterion of reality plays towards the close of antiquity there is a sign of
a weakening of the impulse to inquiry.\footnote{\textsc{Deussen}:
\textit{Die Philosophie der Griechen}. 1911, p.~410.} It is precisely new
and significant thoughts that come forward in the Academic and the Sceptical
school, points of view which Plato and Aristotle in their intellectual
enthusiasm had not allowed to come into their own. The insight that the
decision as to what can be acknowledged as reality is determined partly by
the relation of objects among themselves, partly by the relation between the
objects and the apprehending and thinking subject, has passed over from the
Academic and Sceptical school to the philosophy and natural science of the
more recent age\footnote{\textit{Den menneskelige Tanke}, p.~104 f.}, as
will show itself in what follows. In antiquity itself the discussion of this
question formed the introduction to a purely empirical science, particularly
in the schools of medicine; and a scientific movement in an empirical
direction was perhaps, as already remarked, contributory to the very origin
of the question in principle. ---

In more recent times it is above all criticism and positivism that have laid
weight on the concept of relation. Kant, whose doctrine of the categories is
the most significant since Aristotle, came to his philosophy through the
insight that the category of relation is the most important of all, in that
our knowledge everywhere aims at connecting objects through the
determination of their mutual relations. It is two species of relation that
become especially significant, relations of dependence and relations of
likeness, and it is in the work of thinking such relations that knowledge
reveals itself as synthesis.\footnote{\textit{Om Kontinuiteten i Kants
filosofiske Udviklingsgang} (Vid.\ Selsk.\ Skr.\ 1893) p.~53.} In Kant too
the contention between the two points of view regarding the concept of
relation asserts itself. He has a certain inclination to regard it as an
imperfection in our knowledge that we place everything in relations, and he
therefore understands a ``thing in itself'' as lying outside all relations.
Some Kantians, especially William Hamilton, have stressed this still more
strongly and have set it up as a thoroughgoing limitation of our knowledge
that work must always be done by setting one object in relation to another.
In a similar way Herbart declared that our relation-setting can give only an
accidental apprehension of things, and that true reality does not stand in
relations at all. But in Kant the essential point of view was nonetheless
that what is characteristic of all science is the relation between objects
expressed in definite laws. And from this point of view the necessity of
relation for knowledge becomes precisely a source of fullness and richness.
--- The more relations that can be found and determined, the greater the
unity and the greater the manifoldness there will be at one and the same
time in our knowledge.

For Kant's speculative successors, relation became more and more a barrier
which one strove to get beyond by seeking a highest, all-embracing unity,
within which all relations, particularly all oppositions, were to stand as
one-sided apprehensions. Fichte stands here as a characteristic transitional
type. He has a clear eye for the fact that when we seek to find and to
determine relations, it is a spiritual activity that displays itself.
Without spiritual activity, no relation-setting: ``Die reine Tätigkeit ist
Bedingung des Beziehens,'' (Wissenschaftslehre\textsuperscript{2}, p.~244).
Relation is nothing received or imposed, but presupposes an act of thought.
And when by such an act a barrier is set, this means only that new activity
is necessary. Spiritual work, the theoretical as well as the practical,
seeks to reach beyond barriers, and complete freedom from barriers stands as
the exalted ideal that can never be realized in any definite state. In
Fichte's later philosophy, however, his problem became another. Instead of
asking how we, through the setting of relations and through struggle with
the barriers thus set, can attain an ever richer truth, he now asks, in
Neoplatonic fashion, how there can be any striving after truth at all, when
at the basis of everything there lies an absolute being.

Hegel's philosophy proceeded, in its historical orientation, from a
standpoint on which human thoughts rested securely side by side, without any
need asserting itself to determine their mutual relations. The arising of
such a need is in his view a phenomenon of dissolution. It is an unhappy
consciousness in which the mutual harmony of thoughts is broken. Then there
is need of a strict spiritual work, in order that the individual thoughts,
in spite of, or precisely by, their mutual oppositions, may come to stand as
members of a great whole. This happens by means of the dialectical method,
which at bottom means a systematization of concepts of relation. The
individual thought leads beyond itself to other thoughts through the
recognition of its limitation. The contradictions that arise through the
isolation of thoughts hereby fall away. --- A critical investigation now
shows, however, that it is at bottom Hegel himself who isolates the concepts,
in order afterwards so victoriously to demonstrate their higher unity. Thus
at the very beginning of his Logic\footnote{\textit{Wissenschaft der Logik}.
1812. I, p.~22---26. --- More recent Hegelians take a critical attitude here
towards the master. See \textsc{B. Croce}: \textit{Cio che è vivo e cio che è
morto della filosofia di Hegel} (1907), p.~21---23. \textsc{Mc. Taggart}:
\textit{A Commentary on Hegels Logic} (1910), p.~17.} he sets ``being'' as
something on its own and ``nothing'' as something on its own, in order then
to show that they cannot do without one another. Every being has, after all,
its limit, beyond which there is nothing, if no new form of being can be
found. And the concept of nothing means nothing other than precisely such a
limit of being. It is the experience of limited being that Hegel resolves
into two concepts standing in contradiction to one another. --- But although
Hegel himself thus often creates the contradictions his dialectic is to
overcome, and although it also goes far too easily for him to pass from
relation to relation, his philosophy is nonetheless a great expression of
the truth that one can think existence only by working one's way forward
from relation to relation, and an expression of the fact that in this very
work the highest that can be thought is present as a driving force.

With full consciousness the French Neo-Kantian Charles Renouvier has placed
the concept of relation at the head of his doctrine of the categories. That
everything is relative --- he says --- this strong word of scepticism was
the last word of critical philosophy in antiquity, and it must be the first
word of the modern method. And he finds a whole method contained in this
category (la méthode des relations): the point is to determine our objects
through as many relations as possible. --- And criticism here meets with
positivism. For Comte the concept of relation lies contained in the very
concept of positive science, which everywhere seeks to put definite
relations, expressed in laws, in place of the absolute but indefinite
essences or forces of theology and metaphysics. From this side Comte
expressly regards Kant as his precursor, just as he also meets with him in
this, that the work of knowledge is conceived as consisting in connecting
objects (tout se réduit toujours à lier) --- partly through classification
(par similitude), partly through derivation (par
filiation).\footnote{\textit{Discours sur l'esprit positif} (1844), p.~20.}
As Dauriac\footnote{\textit{Contingence et catégorie} (Revue de Métaphysique
et de Morale. 1916), p.~499.} has shown, Comte's classification of the
sciences is laid out on the same principle as that which later lies at the
basis of Renouvier's ordering of the categories, namely in such a way that
one proceeds from the simple and fundamental to the complex and derived.

Goblot\footnote{\textit{Traité de Logique} (1918), p.~138.} objects to every
attempt to set up a doctrine of the categories that our knowledge is still
too provisional, both as regards the final conclusions and as regards the
fundamental principles, for a system of fundamental concepts to be capable
of being set up. This objection does not touch Renouvier, who is clear that
no rational proof can be given for the validity of the categories or the
completeness of the list of categories. They are found by analysis of
experience, are set up tentatively, and must always be confirmed anew by
experience.\footnote{See Séailles: \textit{La philosophie de Renouvier},
p.~91 f.} Renouvier therefore reproaches Kant with having attempted a
deductive grounding of the doctrine of the categories. In peculiar
opposition to this, Renouvier's disciple Hamelin\footnote{\textit{Essai sur
les éléments principaux de la représentation}. (1907).}, all the while
holding fast to his predecessor's table of categories, has sought to give a
dialectical derivation by showing that every category leads over into its
opposite. But he confuses, as Hegel so often does, the correlation of
concepts with a relation of mutual contradiction or opposition between them;
such a relation comes forward only when the individual concepts are taken
out of their correlation through an isolating abstraction, which is then
cancelled again and thereby produces a semblance of dialectic. Hamelin's
dialectical transitions are accordingly very artificial. --- Beyond an
analytical method one does not get, as Renouvier has very rightly seen, in
the doctrine of the categories.

The categories are sought and found by analysis of the actual work our
thought performs when definite questions are put to it. Every table of
categories can therefore have only a provisional character, and the more
special categories cannot be derived from the more general and
comprehensive, although on closer investigation they may show themselves to
be more special forms of them. As everywhere where we have to do with
concepts drawn from experience (in this case from the actual work of
thought), there is an inverse relation between content and extension. The
more elements a concept contains, the narrower the circle of objects for
which it holds. In the doctrine of the categories this acquires a
particularly important significance through the relation between the
concepts of value and the more universal (fundamental, formal and real)
categories. But, as I have sought to show in the fifth chapter of the
treatise ``Totalitet som Kategori'', since all value stands in relation to a
certain whole and to that whole's conditions of existence, there is
nonetheless an inner connexion between the concepts of value and the other
categories.

It is now the task of the present treatise to throw light on the concept of
relation more closely by seeing it in its significance for the various
philosophical problems. It will hereby be expedient not merely to hold to
the concept of relation itself, but also to investigate the principle that
corresponds to the concept, the principle of relation, which I (``Den
menneskelige Tanke'', p.~255) have formulated thus: Every object shall be
considered in as many relations as possible. Where we can find no relations,
we obtain no knowledge. The relationless is therefore rightly declared by
Kant to lie outside all science and to be only a purely negative concept.
Our thinking works like a pair of compasses: both of its legs must find a
place, and we then obtain knowledge through the relation between these two
placings.

When the treatise is said to deal with ``relation'' and not with
``relativity'', it is because the expression ``relativity'' is popularly apt
to be misunderstood, in that one supposes (as in their way the old Sceptics
did) that when something is relative, it means at bottom that it is invalid
or indifferent, has no objective significance. In opposition to this
misunderstanding it must be maintained that all truth finds its expression
in the demonstration of quite definite relations and holds only within them,
but then also holds absolutely there (provided, that is, that the relations
are rightly determined). That all necessity, as Kant already taught, is
hypothetical, and thus holds in relation to its presuppositions, does not
mean that such necessity should not be real necessity. In all necessity the
presuppositions must be asked after, for necessity consists in a relation
between ground and consequent. Such a relation comes about only through the
activity of thought. And from this point of view too the expression
``relativity'' is unfortunate, in awakening the notion of a resting property
or state. But relation consists in referring, in a definite setting of
thought's two compass-legs. There lies in the concept of relation a
direction to the work of thought.


% ============================================================
%  2. THE CONCEPT OF RELATION IN PSYCHOLOGY
% ============================================================

\section{The Concept of Relation in Psychology}
\label{sec:2}
% pp. 16--26
% Range corrected from the skeleton's "16--25": the closing sentences on
% Hume and Spinoza run on across the page turn, and the section 3 head does
% not open until after them on p. 26.

Every attempt to give even a merely provisional explanation of what is meant
by conscious life or mental life will have to lay particular weight on the
fact that everything that is sensed, presented, felt and willed exhibits
certain definite relations without which it cannot be thought. Attempts have
certainly been made to conceive mental life as a chaos of relationless
elements; but such attempts have not, in spite of the energy that sometimes
lies at the basis of them, been able to lead to an intelligible psychology.
The point is that from its first arising every mental element is woven into
a connexion by which it is determined, and through which it can itself
exercise a determining influence in mental life. It is hereby that mental
life constitutes a whole, so that not least here the above-mentioned
connexion between the concepts of totality and relation asserts itself. In
such phenomena as memory, comparison, wonder and preference, the character
of conscious life plainly displays itself, and in them there is precisely
more than a chaos, in that the various elements which psychological analysis
can distinguish in them are throughout determined in their peculiarity by
their mutual relations. And all experience points to more primitive and
simple mental phenomena having to be conceived on the analogy of these; in
any case language has not formed special designations that could express the
difference. We shall come to see in what follows that this hangs together
with the degree of consciousness that can be present with respect to the
relations of conscious life itself.

It is the thoroughgoing significance of the concept of relation in conscious
life that makes it justifiable, as Leibniz and Kant especially have
impressed, to characterize consciousness as synthesis. A relation can assert
itself only when its terms are comprehended together; the relation is,
indeed, precisely the expression of a connexion, of the terms' belonging to
a whole.

The already mentioned double character of relation, as limiting and as
carrying further, appears plainly in psychology. The peculiarity of every
mental element rests in part on its relation to other mental elements, and
is to that extent limited by this relation. But precisely by being
determined through relations, every element, when it appears suddenly and
without clear connexion with other elements, points beyond itself and gives
rise to a need for supplementing.\footnote{In my treatise on Plato's
  ``Parmenides'' (Vid.\ Selsk.\ filos.\ Meddelelser.\ 1920.\ Kap.\ 3) I have
  drawn attention to the fact that while Plato found it difficult to find a
  logical transition from one concept (an Idea) to another, he entertained as
  a psychologist no doubt that even a suddenly emerging and in the highest
  degree valuable experience had its definite place in the series of
  experiences, conditioned by a whole series of endeavours and experiences
  having gone before in a definite order.}

As regards particulars, I must refer to my presentation of psychology, which
is built precisely on the fundamental thought that the validity of the law
of relation for mental phenomena establishes the character of consciousness
as synthesis. Only a few principal features shall be recalled here.

As regards sense-sensations, it long remained undecided how far the concept
of relation, and with it the concept of synthesis, held for them. In spite
of Hobbes' energetic assertion that to have a single unchanging sensation
would be one with having no sensation at all, one was inclined to conceive
sensations as a chaos which ``the understanding'' was first to order. In the
course of the 19th century one came to see the matter otherwise. Chevreul's
discovery of the effect of contrast showed that the quality of sensations
rested on the relation to the quality of preceding or simultaneous
sensations, and Fechner demonstrated an analogous relation for the intensity
with which the individual sensations assert themselves. --- One has been
inclined to conceive it as an imperfection that we do not have what has been
called an ``absolute'' sensation of anything, and even Helmholtz regarded
successive contrast-effect as an illusion. But no ``normal sensation'' can
be demonstrated whose quality would be the ``right'' one. And in any case
such a normal sensation would itself be conditioned by certain relations,
which one then for one reason or another preferred to other relations. ---
We stand already here, with these elementary phenomena, before something
that will show itself to hold for the highest intellectual functions. ---

If we pass from sensations to presentations (which in their simplest form
are reproduced sensations), it shows itself that the condition for a certain
definite presentation arising at a definite moment is the relation in which
this presentation stands to other presentations. It has indeed been
attempted to ascribe to every presentation a sort of impulse to
self-preservation, so that, as soon as other presentations (or sensations)
do not lay claim to the psychical energy, it would ``of itself'' step forth
into consciousness. But even so there are definite relations --- inhibition
or non-inhibition --- on which the appearance of the presentation depends.
And the very independence with which the individual presentation asserts
itself is conditioned, as experience shows, by its relation to the whole
remaining conscious life, and can be understood only through this relation.
There will always be something unnatural, and finally impossible, in keeping
a presentation isolated from everything else in consciousness. Each single
presentation has a tendency to call forth other presentations according to
definite laws. Here again the synthetic character of mental life reveals
itself. Most psychologists are accordingly agreed that the so-called laws of
association are different forms of whole-association, that is, of the need
for the reproduction of the entire state of consciousness in which the
individual presentation was once a member. ---

In the domain of the life of feeling, the law of relation was discovered
long before its significance for other sides of mental life was recognized.
Most plainly the significance of relation appears in wonder, whether this be
a peculiar feeling or a member of more complicated feelings. But in all
pleasure and displeasure, joy and sorrow, the relation between the present
state and preceding states is of decisive influence. The happiness that is
felt rests not merely on the definite magnitude of the gain or the progress
experienced, but also on the relation to the level of feeling that was
already present beforehand. The more different relations that are decisive
for a feeling, the more peculiar and rich that feeling will be. With every
such richly formed feeling the question will be whether it can maintain
itself in face of the new relations in which mental life may be placed. Thus
for example when the humorist is met by tragic fates\footnote{\textit{Den
  store Humor}, p.~102---104.}, or when a hard-won consolation is to stand
its test in face of new dispensations.\footnote{\textit{Oplevelse og
  Tydning}, p.~150---152.} ---

As little as sensation, presentation or feeling could arise under wholly
uniform inner and outer states, just as little could there be any willing.
To will is to prefer, and thus presupposes that inner or outer differences
assert themselves. Reflex movement and instinct already presuppose changes
in the inner organic state or in the outer relations. Reaction presupposes
action. This holds both for involuntary and for voluntary willing. There is
choice of all degrees, but a relation of difference is always presupposed.
All deliberation, the involuntary as well as the voluntary, aims only at
bringing such relations of difference out as plainly as possible. The
various possibilities of action are measured by their relation to the
momentary or hitherto existing state of mental life. There may here be a
struggle whose outcome rests on the relations to the individual's character
and past. One has sometimes conceived motives as absolutely independent
elements, independent of the will itself. But every motive is itself a
preferring, a willing. What can become a motive for an individual rests on
the individual's character and previous history. Motives are never
relationless. It is of decisive significance at what point of an
individual's development a possibility of action appears; on that will
depend whether a certain motive can come about at all, and whether it can
become the victorious one.

\bigskip
But it is one thing that a proportionality asserts itself everywhere in
consciousness, and another how far consciousness itself becomes aware of the
relations that are decisive in the individual cases. Many psychological
problems arise from one's not discovering the relations, these being perhaps
discoverable only by objective, historical investigation --- an
investigation the individual is not himself in a position to undertake, at
any rate not at the time of the experience itself.

Relations of contrast assert themselves from the first, both as regards
qualities and intensities, without our noticing it. We are taken up with the
life and the strength with which the new element brought out by the contrast
appears. Thus when a colour stands ``saturated'', we readily forget the
relation by which the saturation is conditioned. We turn towards a new
feeling that comes forward with life and fullness after a feeling of the
opposite kind, and forget the vanquished for the victor. An intention or a
decisive resolution may announce itself as ``a blow on the head'', in that we
forget the involuntary or voluntary deliberation that has paved the way for
it through an interplay of memories, feelings and impulses. We do not always
notice the course of association to which a presentation owes its coming
into being; it may then stand as a sudden, unmotivated inspiration.

We have on the whole an inclination to dwell on certain salient points in
our mental life and to overlook the relations that condition the transition
from one such point to another. William James, who has excellently
illuminated this peculiarity of mental life, distinguishes between
resting-places and flying-places within the stream of consciousness, and
holds that we have an inclination to keep to the resting-places. He
expresses himself more plainly when he speaks of substantive and transitive
parts of our inner life, and holds that we are inclined to ascribe to the
former an absolutely independent subsistence, in that the transitions are
not easy to follow and often cannot be directly observed at all, but must be
inferred to have been there.\footnote{\textit{Principles of Psychology}. I,
  p.~243 ff.} Such substantive elements are found precisely in
contrast-determined qualities or intensities, in immediate recognition (the
quality of recognition), and in a suddenly emerging presentation, a marked
feeling, a surprising decision.

James has besides drawn attention to the fact that relations are sometimes
felt as qualities, thus in those mental states which are expressed in
language by ``and'', ``but'', ``in spite of'' and the like. It is a kind of
sensation (or feeling) of connexion or of opposition, which can be just as
immediate as the ordinarily so-called sensations. The so-called quality of
recognition is of the same kind.

Here may also be reckoned the remarkable sensation of direction in our
mental life which we can have. To use the comparison with the compasses, we
may say that while the one compass-leg stands fixed (determined by our past
and its results), we can feel that the other compass-leg is extended in a
certain direction, without our being clear at what definite place it will
come to stand. It is now more an indefinite expectation, now more an
indefinite yearning, that here appears.

On the transition between psychology and epistemology lies the question
whether, and how, a relation of which one has become more or less clearly
conscious can be transferred to terms other than the original ones. An
important way in which this can happen is the following. We involuntarily
order our objects according to likeness and difference. Thereby series are
formed, and it may then show itself that in such a series a term can be
pushed out and the relation nonetheless continue to hold between its former
predecessor in the series and its former successor. In the series ABC,
perhaps B can be pushed out and the relation between A and C be the same as
that between A and B.\footnote{Cf.\ \textit{Den menneskelige Tanke},
  p.~162 f.} It is such a case that formal logic constructs as a relation
between premisses and conclusion. The premisses become, for example: A = B
and B = C, and the conclusion becomes A = C. Relation --- which precisely
for that reason belongs among the fundamental categories --- asserts itself
generally prior to all formation of concepts, judging and inferring. Already
within sense-, memory- and fantasy-intuition, elements are ordered and
reordered on account of their mutual relations, without there yet being any
question of thinking in the narrower sense of the word. I refer here to my
investigations of the relation between intuiting and judging in the first
main section of ``Den menneskelige Tanke'' (especially p.~53---59). Of
particular interest here is the way in which intuition is articulated quite
involuntarily and unconsciously, without a judgment being formed. A striving
asserts itself in conscious life to dwell in intuition as long as possible
without forming concepts and judgments or drawing inferences.

In the series ABC, A remained standing, but C could enter into the same
relation to A as B had had. The question now arises whether one cannot move
both compass-legs, so that they are set at new places, but the relation
between them continues to be the same. In the simplest case such a shift is
present when there is a relation between two new terms (X and Y), and this
relation is then recognized as the same as the one already known between two
earlier terms (A and B). Such recognition can --- like the recognition of
qualities --- be either immediate or mediate; in the latter case it
presupposes explicit measurement, or else there are several relations
serving as intermediate links between the two relations that are identified.
The intuition in which, according to the legend, it dawned on Newton that
the relation between the falling apple and the earth could be the same as
the relation between the planets and the sun, was at first immediate; but it
became mediate through the measurements and calculations that led to the
definitive setting up of his doctrine of gravitation. On the whole many
questions and problems arise precisely from the possibility of identities of
relation appearing in one way or another before consciousness. Both
questions of determination, which are due to lacking relations, and
questions of decision, which are due to conflict between different possible
relations, can lead thinking further. Already within ordinary consciousness
(sens commun) such questions, with attempts at answering them, can appear,
and what we call science rests only on the strict grounding and formulation
of the questions and sharp investigation of whether the answer really is a
solution.

It is an increasing degree of explicit consciousness of relations that leads
over from involuntary thinking to logical and mathematical thinking, which
handles relations as simple qualities that can become terms in new series,
and these again in series of series. An example of the transition is shown
by the development of the concept of number from quality of number through
ordinal number to cardinal number, and thence again to the various concepts
of function.\footnote{\textit{Den menneskelige Tanke}, p.~189---193.}

\bigskip
A definite limit for thinking would be present if there were series whose
terms could be interchanged without further ado, because the relation of
likeness or difference was absolutely the same between all the terms. Such
series would be either chaotic series of difference or absolute series of
identity. Whether they occur in existence cannot be decided in advance. But
from the standpoint of psychology it must be said that they do not occur in
actual mental life, where at most approximations to them can be found. It is
a fundamental law in psychology that the order of the addenda is not
indifferent. Where there seems to be chaos or absolute identity between
psychical elements, it will be the psychologist's task to find, in the first
case, likenesses or definite relations of dependence between the apparently
absolutely different elements, in the second case, differences and the
relations of dependence following from them within the apparently absolutely
alike. Of chaos and of absolute identity one can speak only from definite
points of view, and in relation to them.\footnote{Anf.\ Skr.\ p.~165; 170.}
Thus the principle of relation asserts itself here too. The work of thought
presupposes a constant reciprocal relation or a constant contention between
continuity and discontinuity, likeness and difference; only through such
contention does thought obtain its definite tasks and can begin to work.
Hume's assertion of chaotic elements as what is properly given meant that on
his conception no work of thought could come to begin. Spinoza's assertion
of the eternal attributes as absolute identities presupposed that all work
of thought was done. The two thinkers ended, each in his own way, in
relations that really were not relations at all.


% ============================================================
%  3. THE CONCEPT OF RELATION IN EPISTEMOLOGY
% ============================================================

\section{The Concept of Relation in Epistemology}
\label{sec:3}
% pp. 26--76

\subsection{The Concept of Relation and the Other Fundamental Categories}
\label{sec:3a}
% pp. 26--33
% Range corrected from the skeleton's "26--32": the closing paragraph on
% identity and substitution runs on to the top of p. 33, where the b. head
% then stands.

The analysis that is to lead to the discovery of the forms and
presuppositions of thought must be directed not merely at scientific work in
the stricter sense, but also at the involuntary thinking we all practise in
daily life, that which is characteristic of what one calls ``sound common
sense''. Out of this involuntary thinking, scientific thinking, which
proceeds from definite concepts and seeks exact grounding, has developed
step by step; and in the transition to science there cannot, however great
and striking the contrast may be, be any complete change in the manner of
thinking. Kant saw this when he stated that even ``der gemeine Verstand''
has certain ``a priori cognitions'' (Kritik der reinen
Vernunft\textsuperscript{2} p.~3), and shows it still more definitely in his
doctrine of synthesis as the fundamental form of consciousness, springing
from ``a function without which we should not have knowledge, though we only
seldom become conscious of it'' (Kritik der reinen
Vernunft\textsuperscript{1} p.~78). We have seen above what significance
this concept has for psychology. Like the other fundamental categories, this
concept is one in which psychology and epistemology meet. In the most recent
time Meyerson\footnote{\textit{De l'explication dans les sciences.} 1921. I
p.~IX---XIV.} has strongly and in a very instructive way emphasized that it
is quite the same thinking that is practised in natural science, philosophy
and sound common sense (sens commun). He does not, however, go more closely
into the stages and paths by which practical thinking successively
approaches science. He thus does not emphasize the significance of doubt as
disquieting, in that practical orientation in the world does not always
suffice, or else leads to conflicting possibilities in face of which one
stands perplexed. It is in the school of doubt that thought comes to make
stricter demands upon itself.

If one wished to define the word ``reason'' in such a way that it can fit
pre-scientific as well as scientific thinking, one might say that it is a
sum of categories, particularly of fundamental categories, which must be
applied at every stage. Among these fundamental categories are to be
reckoned the three pairs of concepts synthesis-relation,
continuity-discontinuity and likeness-difference.

Relation is (like synthesis) a fundamental form which constantly recurs in
all special forms of thought. It cannot then be surprising that it asserts
itself particularly within each of the named pairs of concepts. Thus there
is relation between synthesis and relation. These two concepts can, namely,
come into opposition to one another, so that comprehending together and
unity come to prevail at the expense of determination through relations, or,
conversely, the terms of the relation assert themselves so independently
that they do not seem to be members of the same comprehended whole. On the
other hand, the need for synthesis becomes stronger the stronger the mutual
tension of the terms is. When Kant says that we only seldom become conscious
of the function that lies at the basis of synthesis, it is because it is
only when the relation assumes the character of opposition, discord and
tension that the comprehending work can properly come to consciousness.
Herein lies precisely the significance of doubt. And the constant
correlation between synthesis and relation shows itself in this, that doubt
and problems exist only so long as different objects, in spite of ever so
sharp an opposition or contradiction between them, are held together in
explicit relation to one another. A problem falls away when one lets one of
the conflicting terms fall away, or when one weakens the opposition between
them.

There can arise a fear of the manifoldness and sharpness of relations, which
makes the secure repose in involuntary syntheses appear particularly
alluring. Thereby such a reaction is called forth as Hamann, Herder and
Jacobi exercised in face of Kant's critical philosophy, and as James and
Bergson signify in face of the science of the present day. And when Tagore
calls upon Western Europeans to take up Indian ways of thought, there lies
in this precisely a call to disregard all the sharp and tense relations that
have been developed in European civilization, and to return to a relationless
state. But however great and deep the search for unity may be that
distinguishes the Vedanta doctrine, of which Tagore is the modern prophet,
there is no way back. Either the relations that assert themselves are
artificial and unfounded, and then criticism will overthrow them, or else
they are founded in actual relations, and then they can be mastered only by
a spiritual labour. Doubt is mastered only either by its groundlessness
being demonstrated, or by the solution of a problem, or by the demonstration
that here an insoluble boundary question is present. When Descartes infers
precisely from the possibility of doubt to the existence of thought, this
meant a recognition of the close connexion between the categories synthesis
and relation.

The relation between continuity and discontinuity is analogous to that
between synthesis and relation, or is a more special form of it. Synthesis
proceeds the more easily the fewer interruptions occur. If possible, such
interruptions are involuntarily taken up into the intuition (in that this is
more definitely articulated) or into the judgment (in that its terms receive
additions that limit or extend). Where continuity prevails, intuition and
reflection take on the character of a smoothly gliding stream without eddies
and waterfalls. One does not ask, since no arrest or wonder motivates any
question. Or, if one asks, the question runs: Why not? Why not go on as
hitherto, when there is energy enough for the expansion to be continued? But
in life and in science more or less sharp divisions of labour assert
themselves, and thereby resistance and opposition which one cannot get round
without further ado. One must then attempt to find intermediate links
between the terms that have come to stand isolated in relation to one
another. Or one must attempt to demonstrate a continuity lying deeper than
the discontinuity that has forced itself forward, a continuity within which
the leaps or breaks before us can come to stand as derived in relation to
the law or the type the new continuity exhibits. From the purely
psychological side this has already been mentioned above; it is on the whole
the task of psychology to bridge the often so sharp oppositions and
catastrophes that mental life can exhibit to immediate observation.
Particularly of interest here is James' demonstration of ``transitive''
states, in contrast to the ``substantive'' ones. Psychological discontinuity
can on the whole signify only a problem, never denote a solution. Analogously,
in biology the doctrine of mutation, which asserts the sudden arising of new
organic species independently of outer relations, signifies new problems.
The founder of the doctrine of mutation himself, Hugo de Vries, at once
stated that the mutations signify just so many problems. The modern quantum
doctrine in physics, according to which energy is not released in a
continuous way but in leaps\footnote{\textsc{Poincaré}: \textit{L'hypothèse
des Quanta.} (Dernières Pensées. Chap.~6).}, is sought to be more closely
grounded through the demonstration of mathematical laws according to which
the release in leaps stands as necessary. In every case theoretical points
of view are sought for the understanding of the discontinuity. Already when
the leaps can be ordered in a certain series, thought begins its work of
finding a new continuity, even if this is for the time being more formal
than real.

The problems that lie here were already raised in the earliest Greek
philosophy. Parmenides' ``One'' meant at one and the same time continuity
and identity; these were not different for him. Discontinuity, on the other
hand, was asserted by Heraclitus and Democritus. Through Platonism and its
opponents the dispute extends right up into more recent times. Kant regarded
this dispute as a rational necessity. His doctrine of the antinomies aims at
demonstrating here an absolute limit for reason. The great founder of
critical philosophy did not see that discontinuity can never be thought's
last word except in the form of a last problem, and his ``antitheses'',
which assert precisely this point of view, are therefore decidedly right as
against his ``theses'', which assert absolute discontinuities.

Continuity and discontinuity are correlata that reciprocally supplement one
another. They denote different points of view and operations, and the
history of science shows how now the one, now the other of these categories
comes to stand in the first rank, but in such a way that the struggle
between them will always begin anew. No investigator has illuminated this
relation better than Henri Poincaré, especially in the following statement:
``This struggle will last as long as men will cultivate science, as long as
humanity will think, for it springs from two irreconcilable tendencies in
the human mind, two tendencies which it cannot lose without ceasing to
exist, namely the need to comprehend, and we comprehend only the limited,
and the need to intuit, and we can intuit only extension, which is
unlimited. Even if this struggle should not end with the final victory of
one of the combatants, it would not for that reason be fruitless, for in
each new struggle the field of battle is another; each time, therefore, a
step forward is taken and a conquest is made, not for one of the combatants,
but for humanity.''\footnote{\textit{Les conceptions nouvelles de la
matière} (in the collective work ``Le Matérialisme actuel'', 1913), p.~67.
--- Cf.\ also the interesting discussion in the French philosophical society
13 March 1910 (Bulletin de la société française de philosophie. Vol.~X),
especially the statements of Perrin, Brunschwicg and Meyerson
(p.~116---120).} --- To this fine and apt statement I must, however, add the
remark that the relation between intuiting and thinking can also have
another character than that indicated by Poincaré. It is often precisely
intuition that holds fast to limited images, while thought, in virtue of the
recognition of law, sees the limitation of every intuition and the
possibility of new terms in the series of images. Besides, there are several
transitional forms between intuiting and thinking, which have in part been
indicated in the foregoing, and there will in each individual case be
definite objects, given in experience, which decide whether intuiting or
thinking shall have the word.

Where continuity can be maintained, the condition will be that no great
differences in time, place, intensity, magnitude or quality assert
themselves between the objects with which consciousness is occupied. Where
likeness prevails, one passes easily and, at any rate apparently,
continuously from the one object to the other. Likeness and difference are a
pair of concepts exhibiting an analogy with the pair continuity and
discontinuity. But while it can be said of continuity that it is not noticed
until it is past, likeness has a somewhat more conscious character. It does
not stand in the close relation to intuiting that continuity does. --- The
greater the likeness, the better will continuity be able to be preserved. It
is only through explicit, perhaps voluntary comparison that one discovers
the difference between the various degrees or species of likeness. I have
(in my Psychology and in ``Den menneskelige Tanke'' p.~64 f.) distinguished
five species of likeness: identity, congruence, composite likeness, likeness
of quality and likeness of relation (analogy). In this series the element of
difference comes forward more and more, and it is essentially through
increasing difference or decreasing likeness that we pass over from
continuity to discontinuity.

The relations likeness and difference are of such essential significance
that it is justifiable to define thinking as comparison. That the two
relations govern all thinking and all categories was already seen by
Porphyry, when he found it necessary to write an introduction to the
Aristotelian treatise on the categories\footnote{\textit{Isagoge in
categorias} (Scholia in Aristotelem ed.\ Brandis, p.~1).}, in which he
explained the concepts genus, difference, species, individual and property.
It was this ``Introduction'' that came to lie at the basis of the many
discussions in the Middle Ages about the objective validity of general
concepts (the dispute between nominalists and realists). The five concepts
cited form a progressive series as regards differences, thus from the
abstract or general to the concrete or special. For the thinking of the
Middle Ages, which was predominantly classificatory and syllogistic, the
question which of these concepts was to be applied in the individual case,
and with what right it was applied, was of the greatest significance. More
recent science here undertakes a great reduction. For it, it is not the
relation between higher and lower concepts that is decisive, but everything
comes down to whether a concept (an object) can step into the place of
another in all respects or not. It is identity that puts the other species
of likeness in the shade. And precisely by the fundamental character of the
principle of identity as the nerve of all inference, and by the sharp
determination of the concept of identity conditioned thereby, modern science
distinguishes itself from ordinary consciousness, which does not take the
difference between identity and mere likeness so exactly. Parmenides had,
moreover, already seen that the nerve of thought lies here, though it was
first Leibniz who laid the concept of identity at the basis of logic,
because it made substitution possible. It is in this connexion indifferent
whether one defines identity as the greatest possible likeness or as the
least possible difference; both definitions testify to the correlative
relation between the categories likeness and difference. --- But at the
decisive turning point denoted by the dominion of the concept of identity,
we pass over from the fundamental to the formal categories.

\subsection{The Concept of Relation and the Formal Categories}
\label{sec:3b}
% pp. 33--45
%
% STRUCTURAL CORRECTION (2026-08-23). The skeleton set α, β and γ here as
% printed \subsubsection* headings with titles. The printing has no such
% headings: on pp. 33, 40 and 44 the Greek mark stands BARE at the head of an
% ordinary indented paragraph, with no title text at all, and the titles
% ("Rene Kvalitetsrækker" etc.) appear only in the Indhold on p. 109. The
% transcription handles this with \addcontentsline alone, the mark run into
% the paragraph; that treatment is mirrored here. Section 3d is different --
% its marks DO carry printed titles run in with an em dash -- so the skeleton's
% \subsubsection* form is kept there.

\addcontentsline{toc}{subsubsection}{α.~Pure Quality-Series}%
\label{sec:3b-alpha}%
α. Through the working out of the formal categories (identity, analogy,
negation, rationality) it becomes possible to obtain exactly determined
relations.

The concept of identity is formed as a natural continuation of the tendency
which practical thinking already has to form series of objects ordered
according to increasing likenesses or decreasing differences. Its formation
makes possible points of view from which not merely an overview of the world
of relations can be won, but by whose help the one relation may perhaps be
derived from the other. Only at this point do the epistemological and the
psychological-historical ways of consideration come into decided opposition
to one another. Psychologically and historically, at most approximations to
identity assert themselves, and the pure concept of identity is perhaps even
rejected, in order that the differences and changes which observation
exhibits may come into their full right. Epistemologically the relation is
considered from the opposite side: differences and changes stand as breaches
of the identity which is the first presupposition of all science, and which
conditions what becomes a scientific problem.

The fundamental significance of the concept of identity for scientific
thinking rests on the fact that its application is the presupposition of the
two most important functions of thought: classification and proof.

Division of the various objects which observation shows us can become of
decisive significance only when there is a point of view from which the
objects, in spite of all difference, can be regarded as identical. In
concepts of genus and species such points of view are given. The individual
objects for which such concepts hold can --- seen from the point of view of
the concept in question --- be regarded as one and the same object, which
displays itself in the fact that any one of them can be used as an example
of the concept. The differences between the individual objects do not
thereby fall away, but can, precisely by being united under the same point
of view, come forward as sharp oppositions. There is perhaps not even a
single one of those properties which the concept states to be ``common''
that occurs in quite the same way in the individual objects exemplifying the
concept. But an analogy nonetheless prevails between the various objects, in
that the relation between their properties is throughout the same.

Aristotle laid it down with great emphasis, both that science concerns only
the general, the content of concepts of genus and species, and that what
really exists is always individual. Of the individual one could in his
opinion form no concept. He has thereby, without being clearly conscious of
it, posed a great problem which extends under various forms through the
whole later history of scientific thinking. Considering the question here
provisionally only from the point of view of classification, it must be
pointed out that concepts can nonetheless be formed whose content is those
properties or relations which, in spite of all differences in an individual
object's several states, can express the object's identity with itself. Such
typical individual concepts there will be possibility and reason to form
more in the domain of the human sciences than in that of natural science.
Between the various states of the same object there prevails, as between the
various examples of the same general concept, only analogy, not identity.

Where neither analogy nor identity can any longer be demonstrated, we stand
before a chaos. But the concept of chaos itself acquires real significance
only through the relation of opposition to identity. There is a difference
between the chaos that stands before ordinary consciousness when it is
overwhelmed by the differences, and the chaos that is posited by a thought
which in serious labour has searched for analogies and identities, and now
at last lays down the work. And even in this latter case it can never be
maintained with certainty that we stand before a chaotic series of
difference. Strictly speaking it is not justifiable to say anything more
than that it has hitherto not proved possible to find points of view of
unity for the objects that seem to form such series. Just as identity can be
defined by the greatest possible likeness or the least possible difference,
so chaos can be defined as the greatest possible difference or the least
possible likeness. The likeness is so small that it has hitherto not been
noticed, just as with the objects we regard as identical we can say that the
differences are so small that they are not noticed.

The concept of negation emerges here as an expression of the fact that an
interruption of the work of thought --- here of the dividing work of thought
--- shows itself necessary, since no further points of likeness can be
found. The history of classifications shows that in such cases one often
forms negative concepts (or concepts containing only negative marks), as
when Lamarck divided animals into vertebrates and invertebrates; later
Cuvier found positive marks for the invertebrates and the possibility of a
positive division of them.

The chaotic series of difference as well as the absolute series of identity
denote two ideal boundary lines for knowledge, boundary lines which can
never be demonstrated in actual knowledge. The one of these boundaries would
denote the impossibility of knowledge\footnote{Stuart Mill attempted to lead
all inference back to an associative transition from one particular to a
particular different from it. But such a transition is, even purely
psychologically, possible only when the first particular is recognized; only
then can it be understood how it can call forth the second particular with
which it was formerly conjoined. It is, however, rather an inference from
analogy that makes it possible to infer ``from particular to particular'':
as a in an earlier observation stood to b, so a' will now stand to b'. But
then there must after all be likeness between a and a'. Cf.\ \textit{Den
nyere Filosofis Historie} (First ed.\ II, p.~377; Second ed.\ II, p.~414).
(In the third edition: III, p.~119). Cf.\ in the most recent time:
\textsc{Ettore Galli}: \textit{Nel regno del cognoscere.} (1919)
p.~159---167.}, the other its conclusion. Between the two lies a scale of
series of thought (some of the most important of which are presented in
``Den menneskelige Tanke'' p.~162---174), denoting different stages by which
knowledge works its way towards its ideal. ---

So long as knowledge aims only at a classification, it lets the differences
subsist and merely seeks to attain an ordered overview of them. In proof, on
the other hand, the concept of identity acquires active or productive
significance, affords a means of demonstrating new truths, perhaps of
finding new objects without constantly relying on observation. Here the
category of rationality emerges, which appears in the relation between
ground and consequent. Its presupposition is an identity between two objects
justifying us in everywhere putting the one in the place of the other. The
two premisses from which the inference proceeds must contain a concept that
is identical in them both. Classification already makes inference possible,
after identical concepts of genus and species (or typical individual
concepts) have been formed. But science seeks to go a step further. It forms
not merely ordered series of known terms, which can be continued only by the
help of new observations, but also series whose continuation can take place
according to the same law by which it began. Here the difference between
logic and mathematics appears. The mathematician sets up not merely the
principle of identity as holding for the units with which he works, thus 1 =
1, but he sets up also the presupposition that out of these units he can
form a series arising by unit being added to unit in always the same way,
thus 1 + 1 + 1 . . . A presupposition such as this last is unknown to pure
logic. For it, repetition of an object has no significance for thought,
however much psychological or pedagogical significance such repetition may
have. It holds to the principle of identity (and the logical principles more
or less directly connected with it). The fundamental mathematical
presupposition is, on the other hand (alongside the principle of identity),
the positive and productive significance of repetition. It is this
presupposition that Poincaré has called loi de recurrence, and in which he
finds the significance of what Kant has called an a priori synthetic
judgment.\footnote{Cf.\ on this \textit{Totalitet som Kategori} (Vid.\
Selsk.\ Skr.\ 1917)}

While classification left the qualitative differences as they were, and
merely ordered them, modern natural science seeks to work up or purge the
differences of quality in such a way that out of them series of a kind
similar to those just mentioned can be formed, thus purely formal series in
which new terms can always be constructed according to the same law which
has already given the preceding terms their place. There are four species of
qualitative difference which have particularly made such a purging possible,
namely time, number, degree and place. To immediate observation they appear
not as ``pure'', but as qualities. The difference between times, between
places, between numbers and between degrees is just as much qualitative as
the difference between yellow and blue, hardness and softness. The history
of science shows what work of thought has been necessary before the four
named series of quality could be transformed into series in which each
single term is fully determined by its place in the series, a place itself
determined by the relation to the other terms. What fills the moments or
happens at the different places, what is counted or measured in degrees, is
now indifferent. Each single term in the series is regarded as though it had
the same content as the others, so that only the place makes a difference.
In brief outline I have described this process of purging in ``Den
menneskelige Tanke'' (p.~186---201), and I shall not go more closely into it
here. I shall only add that the four processes of purging do not take place
independently of one another, but reciprocally support one another, and that
this becomes possible through the categories arising sporadically in
general; a certain development of the one category can thereby become
significant for the development of the other. Thus the concept of time has
developed with the support of the concept of number and the concept of
space. And it seems as though the concept of number has on the whole been of
great help in the purging of the three other categories. The process by
which unit is added to unit is simpler and easier than that by which moment
is added to moment, degree to degree and place to place. Kant already saw
that in counting there was the simplest example of ``Synthesis des
Gleichartigen''. Comte has emphasized this still more strongly, and Gauss
too has stated that number is in a far higher degree than space a product of
our own mind.%
\footnote{In the ``Reflexionen'' edited by Erdmann, \textsc{Kant} says that
it is really only numbers and magnitudes that can be constructed out of pure
reason, because with them mere repetition is needed. (Cf.\ \textit{Totalitet
som Kategori}, p.~36). --- \textsc{Comte} found in pure speculations about
number ``la véritable origine de tout le système scientifique''
(\textit{Discours sur l'esprit positif.} 1844, p.~100. --- \textsc{Gasus}
declared in a letter that number is a product of our mind, while space has
also a reality outside it. (See the citations in E. \textsc{Mach}:
\textit{Erkenntnis und Irrtum}\textsuperscript{1}. p.~384).}
% The Danish note prints "Gasus" for Gauss, and the parenthesis opened at
% "(Discours" is never closed; both are logged as printer's defects at their
% sites in transcription.tex and are reproduced here as printed.

While typical individual concepts, which concern concrete objects, can be
developed only with constant regard to and borrowing from observation, there
has through the named processes of purging been provided the possibility of
forming series of thought which can be continued according to the same
principle that lies at the basis of their formation, and from which demands
can be made of the objects observation exhibits.

\addcontentsline{toc}{subsubsection}{β.~Identity and Relation}%
\label{sec:3b-beta}%
β.~In the formation of the mentioned pure series of quality the concept of
identity everywhere lies at the basis, in that they are characterized by the
concepts time, number, degree and place being freed from all content that
fills the time, constitutes the number, exhibits degree, occupies place;
with the content all real differences within the series fall away. That they
must still be called series of quality comes from the differences between
time, number, degree and place among themselves continuing to subsist. There
is, as there will later be occasion to mention, a tendency to let number
become dominant at last and to make all knowledge a system of equations.
Yet here, as Meyerson has shown\footnote{\textit{De l'Explication dans les
Sciences.} I p.~264; II p.~167; 204---220.}, algebra and geometry contend
for precedence. The question itself belongs under the real categories.

But however high a degree of exactness the knowledge working with the named
series may attain, the concept of relation announces itself even in these
pure regions. This shows itself especially at three points.

Identity is itself a relation, an object's or a thought's relation to itself
under different circumstances. Whatever definition of identity one employs,
determination through relation is not avoided. Whether one speaks of the
greatest possible likeness or the least possible difference, the tinge of
relation does not disappear. And this tinge becomes particularly strong when
one employs Leibniz's experimental definition: identical are those thoughts
or objects which without cancellation of validity (salva veritate) can
everywhere be put in the place of one another.%
\footnote{As mentioned in ``Totalitet som Kategori'' (p.~41 f.), various
conceptions have asserted themselves among Danish mathematicians in face of
the question whether their science requires fundamental propositions
expressing absolute identity, or whether approximation to such is
sufficient. This opposition recalls the two different definitions of
identity --- as greatest possible likeness or as least possible difference.
In the last work we have from \textsc{Zeuthen}'s hand, he states
(\textit{Hvorledes Matematiken i Tiden fra Platon til Euklid blev rationel
Videnskab.} Vid. Selsk. Skr. 1917, p.~80) that the applicability of Euclid's
geometrical system rests on \emph{that} the empirical lines have the
properties stated in the postulates, which Euclid presupposes, or if one
will, it depends on \emph{how far} they have them. By this ``how far'' there
is surely an allusion to the conception which \textsc{Hjelmslev} has
asserted, and according to which ideal definitions such as the Euclidean are
not needed, but one can proceed from lines and figures that can be drawn.}
By making everything depend on the possibility of substitution in all cases,
one plainly enough employs precisely the relation between the thoughts
(objects) in question and others; validity is preserved only if such
relation is not altered when the substitution takes place. Even with the
identity 1 = 1 the question may arise whether we really dispose of units
that can everywhere be put in the place of one another. In the discussion of
Fechner's attempt to give Weber's law a strictly mathematical formulation
this question came forward.

But secondly, the setting up and application of the concept of identity
always presupposes a definite point of view from which alone it holds. In
classification it is plain enough that the identity holds only for certain
marks which constantly recur, so that there is at least analogy between
their forms in the individual cases. But the identity which is the nerve of
every proof also holds only from a certain point of view. Identity is always
attained through purging, that is, by disregarding points of view that do
not concern the matter under discussion. When the various forces of nature
are determined as energies, that is, measured by the capacity to perform
work, there is thereby provided a point of view of unity, just as when
otherwise highly different things are considered solely and alone under the
point of view of weight. Right up to the appearance of spectral analysis one
regarded, in accordance with Newton's physics, the globes in space
essentially under the point of view of weight. It is by adopting a purely
quantitative point of view --- a point of view which again becomes possible
only by one's holding to certain definite properties --- that identity can
be provided between otherwise highly heterogeneous objects. Just as in the
descriptive sciences quality can be reduced to identity by means of concepts
of genus and species, so in the exact sciences quality can be reduced to
identity by means of quantity.\footnote{Cf.\ \textit{Den menneskelige
Tanke}, p.~177---180.} Plato has given profound hints in both directions; as
regards the last point, Galileo took up his thought in an ingenious way.

In Greek thinking a kind of cult was often practised towards the ``One'',
the sum of the pure identities. It appears in Parmenides and Plato and
reaches its high point in Neoplatonism. For the thinking of the more recent
age, identity is not an object of contemplation, but is a great means of
thought, as is aptly expressed in what I have above called Leibniz's
experimental definition of identity. And not less essential is it that we go
a step further back and ask from what point of view the identities hold. We
do not on that account miss a world of Ideas. On the contrary, the points of
view from which identities become possible are for us the true world of
Ideas; and it is a consideration which epistemology contributes to the
discussion of the problem of existence (the cosmological problem), that the
fact that there are such points of view, by which the strictest knowledge
men can attain is possible, must necessarily become of great significance
for a final, rounding-off world-view. ---

Thirdly, it is precisely the formation of the qualitative series of identity
that makes the relations come forward more sharply than is the case before
the purging. In the named series each term is determined exclusively by its
place in relation to the other terms. A moment acquires, quite apart from
its content, an individual character by following after certain moments and
preceding certain other moments. The moments serve reciprocally to measure
out one another's individuality. The same holds of number: each number in
the series is an individuality, precisely in expressing a definite relation
to other numbers. Malebranche already remarks that whole numbers are
relations just as much as fractions are, although one does not always notice
it, because whole numbers can be expressed by a single sign. As Gauss later
stated: ``Every real whole number represents the relation of an arbitrarily
chosen initial term to a definite number in the
series''.%
\footnote{\textsc{Malebranche}: \textit{Recherche de la vérité.} VI, 1, 5.
--- \textsc{Gauss} cited in \textsc{Cassirer}: \textit{Substanz und
Funktionsbegriff}. p.~72. --- This is, however, expressly denied by
\textsc{Bertrand Russell} (\textit{Introduction to Mathematical
Philosophy.} p.~64), for whom the number m is no relation; it is otherwise
with $+$~m and $\div$~m. A number is for Russel's Platonizing conception an
existence in and for itself, apart from the relations in which it is seen.
Therefore Russel also distinguishes definitely between $\frac{m}{1}$, which
is a relation, and m itself.}
% The Danish note prints "henægtes" for "benægtes", and "Russel"/"Russel's"
% with one l against "Bertrand Russell" earlier in the same note; both are
% logged as printer's readings in transcription.tex and are kept. The spelling
% cannot be carried into English, so the sense is translated and the surname
% left as printed at each site.
The cult of certain numbers which is so frequent in popular religions became
possible only by numbers being taken out of the series in which alone they
have their individuality. --- What holds of moments and numbers holds
analogously of degrees and places. This needs no closer demonstration in
this connexion. As regards place (space), another, more special connexion
will give occasion to return to the question. Of interest, however,
precisely here is Poincaré's remark\footnote{\textit{Science et Méthode},
p.~113.} that the relativity of space and its uniformity are one and the
same thing. ---

\addcontentsline{toc}{subsubsection}{γ.~The Relations of Negation and
Rationality}%
\label{sec:3b-gamma}%
γ.~To the formal categories we have besides identity and analogy also
reckoned negation and rationality. Nor in their case is any detailed
demonstration needed that relation constantly lies at the basis. --- Every
No presupposes a (possible or actual) Yes. A negation is purely logically a
sign of a path that cannot be travelled, and thus presupposes attempts to go
such a path. Or it presupposes a comparison between two intuitive images, of
which the one has elements lacking in the other. Mathematically, negation
can signify a call to form a series in the opposite direction from the
earlier one. --- Rationality signifies in itself a relation, a relation
between premisses and conclusion. And as far as the demand for grounding
goes, so far must the presuppositions of our judgments be asked after. An
unconditional necessity does not exist. Necessity expresses only the
relation between ground and consequent, a relation that must constantly
recur so long as there is thinking, whether it be possible to find it in the
individual case or not. Just as in classification we came upon the concept
of analogy, in that, where concepts of species can no longer be found, we
stop at analogies between the individual objects (in so far as such
analogies can still be found at all), so, as regards proof, we stop at the
so-called inference from analogy. Aristotle has aptly called it inference
from an example. From the relations within one object we are led to surmise
corresponding relations within another object. The more examples we hereby
have to proceed from, the more does the inference from analogy approach
induction. But it always acquires (like induction, for that matter)
significance only as a method of discovery, that is, as an occasion for
setting up propositions which can thereafter be more closely tested; it can
never become proof. Therefore it has aptly been called a good servant but a
bad master. That analogy too always rests on definite presuppositions needs
no closer demonstration.

\subsection{The Concept of Relation and the Real Categories}
\label{sec:3c}
% pp. 45--57
% As in section 3b, the α and β marks stand BARE in the printing (pp. 45 and
% 55) with no title text; the titles are the Indhold's and go to the ToC only.

\addcontentsline{toc}{subsubsection}{α.~Causality, Reality, and Relation}%
\label{sec:3c-alpha}%
α.~The question naturally arises whether it is possible to apply the series
of which the formal categories show the possibility to objects given to us
quite apart from our own constructions; and that again is to ask whether out
of such objects series can be formed which in some degree or in some way
correspond to those formal series and can thereby share in the necessity
distinguishing the latter --- thus whether a rationalization of given
objects is possible. The relation between ground and consequent is, after
all, the model of all necessity.

This question gives occasion for a discussion of the real categories, namely
causality, totality and development, which I have presented in ``Den
menneskelige Tanke'' p.~207---237, and here return to only in order to see
what part the concept of relation plays in their domain.

The real categories have historically not waited for the formal ones. Rather
it is they that have given the impulse to the working out of the latter. The
development of science out of practical life and sound common sense appears
strikingly here. The real categories, particularly the concept of cause,
have their actual origin in that need for orientation in existence which is
a condition of life for animals and for men. The development of thought out
of doubt shows itself strikingly here. When disappointment, resistance or
contradiction puts an end to the involuntary trust in all emerging
sensations and presentations, there is no other way to go than to seek new
sensations and presentations which can exhibit a firmer connexion than the
one that was broken, and which can make the secure continuation of life
possible. But this means that the circle of relations is widened until the
various objects show themselves to stand in fixed, unavoidable relations of
connexion to one another. It is causality or lawful connexion that lies at
the basis of all conviction of reality which is possible after the
involuntary trust in the immediately given connexion is broken. This
practical criterion of reality is deepened and made precise in science
through closer testing and determination of the relations conditioning the
connexion between the objects. Beyond all relations thought does not get.
Of the reality of something that should stand in no relation at all, no
grounded conviction can be possible. Reality and relation belong, after all,
together. Every establishment that something is real consists in the
demonstration of a fixed relation between this something and the circle of
objects which hitherto stand as real through mutual fixed connexion; the one
compass-leg of thought stands in this circle, and with the other compass-leg
the relation between this circle and that whose reality is in question is
measured out. If doubt is raised about the reality to which one has hitherto
held, then the first compass-leg too must be moved, and so on. There will
constantly be use for both compass-legs in the knowledge of reality. If one
is forced to stand on the one without being able to find a secure place for
the other, one is in a state of doubt, expectation or longing.%
\footnote{Thomas Aquinas attempted to free the theological concept of God
from all relation by maintaining that the world indeed stands in relation to
God, but God not to the world. (Cf.\ on this my \textit{Religionsfilosofi}
§~22). Thomas here attached himself to mysticism, particularly to Hugo a St.
Victore, who again is influenced by Neoplatonism. Proceeding from this train
of thought, a modern Catholic theologian has criticized Spinoza for not
having dared to take the step of assuming that to the world's dependence on
God there corresponds no dependence of God on the world.
(\textsc{Dunin-Borkovski}: \textit{Der junge Spinoza}, p.~352). But a
relation always presupposes two terms, and however one may stand towards
Spinoza's concept of God, one must acknowledge the consistency with which he
holds fast to the concept of relation on this point. Cf.\ my treatise
\textit{Spinoza's Ethica}, p.~26; 32---35; but also p.~96 f).}
% The Danish note closes with an unopened closing parenthesis, logged as a
% printer's defect at its site in transcription.tex; reproduced here as
% printed.

On account of the close connexion between causality and the criterion of
reality, which shows itself in our having grounded conviction of reality
only where we can connect objects as cause and effect, it was a
revolutionary step that Hume took in doubting the right to apply the concept
of cause. He shook both the scientific knowledge of existence and the
practical criterion of reality. But his criticism rested on the
presupposition that cause and effect were two wholly different things ---
terms in a chaotic series of difference. Here precisely the fixed connexion
is lacking, since one can place the terms of such a series in any order
whatever. Now it is a fact that we nonetheless sometimes stand before series
whose terms cannot be interchanged, and it was in this fact that Kant took
his point of departure as the basis at once for the principle of causality
and for the concept of reality. And the history of science shows how the
external relation between cause and effect in many domains gives place, in
the course of investigation, more and more to a relation of continuity
between the objects standing in causal relation to one another, and this
again perhaps to a relation of rationality, in that later terms in the
series of events can be derived from the law of the sequence of the
preceding terms.

However continuous and rational the series formed in this way may become,
the concept of reality nonetheless continues to stand as an ideal which can
never be fully covered by observation. New intermediate links will always be
findable, and new objects may emerge whose continuous and rational relation
to the earlier ones is to be demonstrated. Reality constantly means a
relation to continuous and rational series formed hitherto. And no guarantee
can ever be given that these series may not have to be replaced by quite
other series, or, to keep to an image used above, that both compass-legs may
not have to be moved.

Herbert Bradley has rightly employed the concept of experience in the for us
inconcludable characterization of existence. Experience is for Bradley the
measure of reality; but since experience is constantly won by the
demonstration of relations, and new relations can always assert themselves,
a complete satisfaction of that measure's demand is impossible. Concluded
harmony and relation are irreconcilable concepts. Bradley is at one and the
same time convinced that we always think in relations, and that a complete
solution of the problem of reality is impossible from the concept of
relation (the relational point of view).%
\footnote{\textit{Appearance and Reality}\textsuperscript{1}, p.~171 (When
thought begins to be more than relational, it ceases to be mere thinking);
cfr. p.~445; 519.}

The solution can, and in this Bradley remains right, never be fully
attained. But the work of knowledge constantly aims at finding new relations
and seeking to bring them into connexion with relations found earlier,
perhaps using them to correct the latter. Without here entering upon what
belongs under the cosmological problem, it must be maintained that the work
of knowledge is itself a reality, a member of the great reality which is
perhaps far more comprehensive than what any work of thought can span. And
when that part of reality which the work of knowledge constitutes can never
be concluded, one might perhaps draw from this the surmise that existence on
the whole is not finished and cannot become finished. Such a
surmise\footnote{Such a surmise I already stated in my University Programme
of 1902 (\textit{Filosofiske Problemer}) p.~61 and shortly afterwards in
lectures in America (Mindre Arbejder. Tredie Række: \textit{En filosofisk
Bekendelse}, p.~26 f.).} presupposes, to be sure, that the form of time
holds for existence, not merely for our thoughts, although at all times,
also now, there are investigators who do not assume this. As regards
Bradley, he definitely denies that thought is a part of reality: The process
which moves within Reality, is not Reality itself. Philosophy is itself but
appearance.\footnote{\textit{Appearance and Reality}\textsuperscript{1},
p.~448; 454.}

Bradley's profound principal work is cosmology (``metaphysics''), not
properly epistemology. He sets reality in opposition to truth: when truth
becomes complete, it ceases to be truth, for then reality enters in its
purity. From this standpoint it follows that although he lays such great
weight on the concept of relation that he defines thinking as finding
relations, he nonetheless conceives thought's search for relations as an
imperfection, because conclusion and totality are thereby excluded. And with
this it again hangs together that he misjudges the significance of the
relations with which the special sciences each work in their own domain. He
has clearly seen the significance of the concept of relation as basis and
measure for all search after truth. But his own system is at bottom a
crystallized ideal of truth, an attempt to think experience as complete
(experience entire, containing all elements in harmony, p.~172), --- an
attempt which cannot possibly be carried through, precisely when one
conceives thinking, as Bradley does, as relational.%
\footnote{In the same direction as the criticism of Bradley given here go
\textsc{Haldane}'s critical remarks (\textit{The Reign of Relativity}, 1921,
p.~203---211) towards this thinker, to whom he stands in other respects so
near. Haldane's work first came into my hands after the present treatise was
completed.} ---

\bigskip
The demonstration of continuous series of objects with non-interchangeable
terms is not yet science's highest ideal. It provisionally remedies the
external relation between what one calls cause and what one calls effect,
which called forth Hume's scepticism; but it does not completely put an end
to it. A step further is led by the demonstration of two-sided equivalence
between different objects. There may then be the possibility both of
deriving the one object from the other and this other from the first. Where
the relation of equivalence is formulated as an equation, this appears
plainly. It is (in theory) indifferent from which side of the equals sign
one begins, and the qualitative difference, as well as the temporal
difference between the two objects, then no longer has any significance; an
absolute relation of identity is attained. The qualities (among which time,
succession, may also be reckoned) then stand as mere ``anthropomorphisms'',
subjective objects about which strict science does not trouble itself. One
here goes a step further than the mechanical conception of nature, which
also disregarded qualities in the narrower sense, but still acknowledged
time as an objective form, in that it saw motion in space as the only thing
that existed, although it was ``falsified'' by men's subjective
apprehension.%
\footnote{\textsc{Thomas Hobbes}, who partly by his own path, partly under
the personal influence of Galileo, became one of the first champions of a
strictly mechanical conception of nature, expresses in his poetic
autobiography the insight he had won (motion as the only thing)
thus:\par
\centerline{Et mihi visa qvidem est toto res unica Mundo}
\centerline{Vera, licet multis \emph{falsificata} modis.}}
It is a significant and paradoxical step that is taken further when one
assumes that time can be reversed, so that the sequence of events is at last
indifferent.

In opposition to such a conception, which is often maintained from the
mathematical and natural-scientific side, and which in philosophical form
has been carried through by the Marburg school, it must be maintained that
the world of pure equations in which strict science ends, or towards which
it at any rate aims, does not lose its significance because it is constantly
held fast that a relation between before and after can never completely
resolve into a relation of pure identity. Because there is a two-sided
relation of equivalence between two qualitatively different objects (such as
heat and motion), it is nonetheless not indifferent which object comes first
and which follows after. Something different happens in the two cases, which
will plainly show itself when the series is carried further through new
terms: even if there is a relation of equivalence between A and B, there may
from B be a transition possible to C which would not be possible from A.
Besides, the transition does not in fact take place equally easily in both
directions (as Carnot's theorem shows with respect to the relation between
heat and motion). But in spite of everything, that strictly mathematical
conclusion of the knowledge of nature has its great significance in this,
that to the purely timeless, logical-mathematical relation between the terms
of the equations there corresponds the temporal transition between
qualitatively different objects in our experience. It is through analogy
that logical-mathematical knowledge acquires its inestimable significance.
In the equations the expert can read off the course of the phenomena, just
as the musician can read off the course of the notes in the score.

This conception was already hinted at by Kant in his doctrine that between
ground and consequent on the one side, cause and effect on the other, there
existed only a relation of analogy, not a relation of identity. It has later
been asserted in a particularly significant way by \textsc{Clerk Maxwell}
(in the second volume of his ``Scientific Papers''). In my University
Programme 1902 (``Filosofiske Problemer'') and in ``Den menneskelige
Tanke'' (1910) I have set forth the same conception, which I have in
``Totalitet som Kategori'' (1917) developed more closely in polemic against
the Marburg school.

That the relation between a series of equations and a series of qualitative
changes is not a logical relation of identity is seen from the fact that
from pure logic and mathematics one could never derive that there is a
temporal world with qualitative differences. It goes with that world of
equations as with Plato's world of Ideas (which at bottom appears here in a
new shape): one can perhaps climb up to it, but the ladder falls away as
soon as one has got up there. And to this comes further that a conclusive
proof of the real validity of the absolute series of identity cannot be
given, since there is always the possibility of new relations whose
analogies with the series of identity found must first be established. Here
again it shows itself that the concept of relation makes truth and reality
ideals which can be attained approximately only through constant labour.

\textsc{Émile Meyerson} has in his epistemological works laid very great
weight on the difference between knowledge of a lawful relation between
objects and a rational explanation of such a relation. Especially in the
recently published great and penetrating work ``De l'Explication dans les
Sciences'' (1921) he seeks to illuminate this difference through a wealth of
examples from natural science and its history, a wealth which everyone who
occupies himself with epistemology must envy him. For my part I agree with
him on the main point, although my terminology is somewhat other than his. I
have already, in the section of my Psychology that deals with the transition
from psychology to epistemology, maintained that science does not, and
cannot, stop at the mere establishment of an unavoidable and
non-interchangeable sequence of two objects, but that the task, after such a
sequence (and thereby a ``law'' in the popular sense of the word) has been
demonstrated, becomes that of finding a continuous transition from the one
term in the causal relation to the other. so that what we call the effect
can come to stand as a continuation of what we call the cause, a
continuation forming an analogy to the way in which in an inference the
conclusion ``follows'' from the premisses.
% The Danish sets a full stop after "det andet" where the sense requires a
% comma (the following "saa" is lower-case and within the same paragraph);
% logged as a printer's defect at its site and reproduced here as printed.
It is the great ideal which the thinkers of the 17th century set up when
they demanded that there must not be more in the effect than there was in
the cause. The continuity found when this ideal is to a certain degree
attained is, so far as I understand, what Meyerson calls rationality, and
what is attained by ``explanation''. In a very instructive example he shows
(I p.~315) how science does not stop at demonstrating a lawful relation
between certain optical phenomena and the chemical structure of a substance,
but that one seeks to find a connected process beginning with the action of
light, which is a certain kind of motion, and continuing through the motions
which this action calls forth in the interior of the molecules, and which
again call forth certain motions in the sense organs affected by the
substance in question. It is according to this scheme a continuous series of
motions that is sought to be established. But in so far as it is possible to
carry this scheme through (and Meyerson very strongly emphasizes the
difficulties in it), it can be done only through the demonstration of a
series of laws (for the reflection of light, for the inner structure of
molecules, and so on), and I do not see what opposition there then finally
remains between ``lawfulness'' (légalité) and ``rationality''. Meyerson
himself accordingly sees clearly that every knowledge of law, often
unconsciously for the discoverer, presupposes an inner bond between the
objects that are tied together in virtue of the law (II p.~289), and he
concedes (II p.~337) that legality and rationality are in reality so closely
(inextricablement) connected that it is only by reflection on the foundation
of science that one can keep the two principles apart. For my part I do not
assume two principles here; it is one and the same principle carried through
with increasing exactness and consistency.

I therefore do not believe either that Meyerson is quite right in his
criticism of Comte, when he formulates this criticism by saying that Comte
indeed acknowledges lawfulness, but not rationality, as the task of science.
Comte himself states definitely (in Discours sur l'esprit positif,
p.~20---21) that the task of science is both explanation and foresight, and
that it ``by its systematic speculations meets our understanding's
involuntary need for unity by carrying through the continuity and uniformity
of our various concepts''. It is another matter that he did not think there
was any prospect of realizing this ideal, partly because the actual
discontinuities seemed to him insuperable, partly because he wished to be a
founder of religion and to make the results of the science won hitherto
(i.e. up to 1830) the dogmas of the new religion. The criticism of Comte
must in my view be formulated thus: he demanded a conclusion of the
knowledge of law and did not see the possibility of carrying it through in
still more special forms and degrees than had hitherto been done. ---

There is yet another point on which I do not agree with the excellent French
thinker from whose works I have learned so much. I do not believe that
science's last goal is the setting up of general propositions, formulated as
exactly as possible. There is yet a work to be done, namely to use the
general propositions, the rational points of view, for the understanding of
the individual existences which subsist only by means of an interplay, a
manifold conjunction of lawfulnesses. Aristotle was wrong that the
individual could not be an object of science. Rather the matter stands thus,
that the immense work science does to find lawfulness and rationality is
itself again the condition for winning an understanding of the individual
wholes which are what actually exists. Only by not forgetting this task is
the eye opened for science's great, infinite goal.

\addcontentsline{toc}{subsubsection}{β.~The Relations of Totality and
Development}%
\label{sec:3c-beta}%
β.~This reminds us that besides causality, totality and development too must
be reckoned among the real categories --- precisely in virtue of the concept
of relation. For the rational points of view that have been found must be
connected with one another and are thereby placed in definite relations to
one another, and through the work on this thought finds new totalities where
formerly there seemed to be only sporadic observations. But every such
totality stands in relation to objects and elements outside itself, and
therefore --- in virtue of the concept of relation --- ever new tasks are
set. When, on the other hand, given objects immediately appear as wholes,
the task becomes not synthetic but analytic; the point is to find the
connexion, the interplay, that makes such wholes possible. What one calls
``things'', immediately given wholes, are understood only through their
properties, and the properties denote just so many relations to other
``things''. Molecules, atoms and electrons are constantly ``things'', wholes
known and understood only through their relations.%
\footnote{A closer exposition of what is here hinted at is given in my
treatise \textit{Totalitet som Kategori.}}

The concept of development comes forward naturally when one lays weight on
causality signifying a continuous process; different stages can then be
distinguished here until a certain conclusion is present, and a
characterization of these different stages can be attempted. But further,
the concept of development hangs closely together with the concept of
totality, in that every whole has its history. --- What in the individual
case is called a whole or a stage of development will depend on the point of
view. What for one consideration stands as a whole stands perhaps for
another consideration as a part; what for one consideration stands as a
stage of development stands perhaps for another consideration as a stage of
dissolution. A closer entering upon the questions arising here would belong
under the discussion of the relation subject---object, or else under the
discussion of the concept of relation in the domain of the concepts of
value, and must be postponed to later sections. ---

Carneades has proved right. We know reality only through the fixed and more
or less transparent relations between the objects observation exhibits. The
progress of modern science in relation to the old Academics and Sceptics
rests on the exactness with which the individual relations are determined,
on the increasing number of relations demonstrated, and on the interplay of
these relations that is discovered. Thereby the concept of reality has step
by step been developed and extended. The Academics did not see the
possibilities of development which the concept of relation offered. They
could not look with humour on the imperfections of knowledge, because they
did not see the significance of a progressive application of the concept of
relation. A peculiar form of humour has, on the other hand, become possible
in more recent times, in that the great ideal of knowledge is held fast
without overvaluation of each step that is taken.%
\footnote{Cf.\ \textit{Den store Humor.} Kap. 6 (Forstaaelse og Humor) og
Kap. 9 § 54 (Humor og kritisk Realisme).}

But to this comes further that behind the Academics' sceptical attitude
there lay at bottom always the popular concept of reality, according to
which reality stands fixed once for all as an absolute order of things with
which our thoughts must agree in order that our knowledge may become true.
It is an impossible relation that is here assumed. Reality cannot itself be
a term in a relation; that would indeed presuppose that we already knew it.
Reality rests, on the contrary, itself on, or makes itself known in, the
relations between our observations. Our thought cannot set its one
compass-leg on itself and the other on something not yet accessible to it at
all. Truth is not unveiling, but bringing forth, and is won only through
constant labour.

\subsection{The Relation Subject---Object}
\label{sec:3d}
% pp. 58--76

\subsubsection*{α.~Introduction.}
\addcontentsline{toc}{subsubsection}{α.~Introduction}%
\label{sec:3d-alpha}%
% pp. 58--62
% Unlike the bare marks of 3b and 3c, section 3d's marks DO carry printed
% titles, run in with the text after an em dash; the transcription sets them
% as headings and drops the run-in dash, and that treatment is followed here.

\noindent The relations hitherto discussed subsist between objects among
themselves and express from various sides the fact that human thought works
by setting one object in relation to other objects and in that way finds the
understanding it can win. The point is at last to lead all objects back to
certain fundamental objects which for the time being stand fixed, even if a
later age's inquiry may seek to reach further back. As \textsc{Ernst Mach}
(Erhaltung der Arbeit. 1872) has expressed it: that one stops precisely at
certain definite objects can be explained only historically. He adduces a
great example (p.~32): ``When we nowadays believe that mechanical facts are
more intelligible than other facts, . . . it is an illusion. It rests on the
fact that the history of mechanics is older and richer than that of physics,
and that we have for a longer time stood on a familiar footing with
mechanical facts. Who dares assert that electrical and thermal phenomena
will not one day stand before us in the same way, when we have come to know
their simplest rules and have become familiar with them? . . . A fundamental
fact is never more intelligible than another''. --- That the place of the
fundamental facts (or, as I call them, the fundamental objects) in our
knowledge can be explained only historically means that it is through the
history of science that we can come to understand why at a given time one
had to stop precisely at certain definite facts and find full explanation by
leading other facts back to them. Mach has in the statement cited foretold
what has really happened. In the twentieth century natural science no longer
regards inertia, weight and the persistence of energy as fundamental facts.
It at any rate sets itself the task of leading them back to other facts. ---
In philosophy one has often wished to rest content with a certain number of
fundamental concepts (categories), and even Kant thought he had here reached
the end. But it shows itself that the fundamental concepts knowledge employs
in its work are as little the same at all times as natural science's
fundamental facts, even though they do exhibit a certain persistent type
through the ages.

It is understandable that investigators in the midst of their energetic work
do not always have at the same time clearly before consciousness the
historically conditioned character of the limits they reach. Meyerson
expresses this thus: science is ontological, only that it puts another
ontology in place of that of sound common sense, whether the new ontology
appears as atomism, as energetics or as pure mathematics.%
\footnote{La science de nos jours est saturée d'ontologie, et les savants,
en dépit de ce qu'ils affirment expressément eux-mêmes, font de la
métaphysique comme ont fait, de tout temps, leurs dévanciers. \textit{De
l'Explication dans les Sciences.} II, p.~174.}
--- It is, however, not necessary that naive ontology should always
characterize scientific work. In the first place, as Meyerson himself
shows\footnote{La science, en menant ses explications jusqu'au bout, finit
par détruire cette ontologie qui, d'abord, paraissait lui être
indispensable. I, p.~50.}, science itself dissolves the ontology which for
the time being seemed necessary. In the second place, the question must
naturally emerge, even for the most enthusiastic investigator, with what
right he now thinks he has attained ``the thing in itself''; that is no
foreign question lying outside science itself, as the history of science
sufficiently testifies. And in the third place the question must inevitably
come forward how, from the ontology one has reached, one can derive
retrogressively the world of qualities and successions from which one
climbed up to what one thought to have attained as the absolute order of
things. This question has the same justification now as it had in antiquity
towards Plato and Democritus. ---

But besides the objective relations through which objects are illuminated
and determined in relation to one another, there constantly asserts itself,
though more hidden, a relation between the knowing subject and the objects
that are its objects. The ancient Sceptics already recognized that here was
a fundamental relation. (See above p.~9). All objective relations are
apprehended and thought through by a human subject and hold, for the time
being, in relation to it. With a change of this subject, its organization
and its conditions, its knowledge too must be able to become another. Here,
as with the objective fundamental objects, there is a historical task,
namely to demonstrate why the knowing subject meets with precisely certain
presuppositions and puts certain questions. This task concerns comparative
psychology and the history of science. Here we have only to do with the
general point of view: without subject no object.

The converse point of view holds too, however: without object no subject.
For we never have a pure subject without objective content and without
determination and conditioning by this content, just as little as we ever
have a pure object, but always (if not with full consciousness) presuppose a
possible observer or thinker. The concepts subject and object are
correlative, like synthesis and relation, continuity and discontinuity,
likeness and difference, and so on. --- Every presentation of a world's
coming into being, from the Mosaic history of creation to modern cosmogonic
attempts, gives a description of how the whole would unfold itself for a
spectator with certain senses and certain presuppositions. The knowing
subject, the spectator, the thinker, has himself a history determined by the
experiences he has had and the way in which he has worked over these
experiences. Modern epistemology is clear that through scientific work there
can develop a subjectivity with the capacity and the need to see objects in
their relations and with an understanding of how it has itself reached the
standpoint from which it sets its tasks and demands answers to its
questions.%
\footnote{On the reciprocal action between the psychological and the
epistemological way of consideration on this point see \textit{Den
menneskelige Tanke}, p.~297---303.}
Such a subjectivity may be called the epistemological subject. Heinrich
Rickert, who has clearly illuminated this point, aptly calls the
epistemological subject a boundary concept in relation to the psychological
subject.%
\footnote{\textit{Grenzen der naturwissenschaftlichen
Begriffsbildung}\textsuperscript{2}, p.~134---138; 308.}
That no doctrine of the categories can be definitive follows precisely from
the fact that this boundary concept can and must always be determined anew,
because it develops in relation to the objective content present at each
time. The world is understood from man's nature and standpoint, but man's
nature and standpoint are again understood from the laws of the world and
the development of the world. Beyond this interplay one does not get, even
if the two points of view may at different times appear with different
weight in scientific discussion.

The weight laid on the relation subject---object leads by no means to
subjectivism, since, as has been said, there is no pure subject that should
be able to bring forth everything out of itself. There is always only an
objectively determined S (an S\textsubscript{o}), just as there is always
only a subjectively determined O (an O\textsubscript{s}). That this relation
does not come to consciousness without further ado is a consequence of the
involuntary trust with which we from the first meet every sensation, every
memory, every fantasy, every thought. On this is grounded the thought which
the Greeks found it so hard to get away from, and which still haunts even
clear heads: what I know must after all be something in itself; were it
nothing, there would be no knowledge either! --- Nor is there at first any
reason for doubt about the very forms of thought in which we discuss our
problems. And for the time being we do not think that we are ourselves a
compass-leg standing at a definite place, from which the other compass-leg
can be placed at various places, but always in relation to the first
compass-leg. We do not at once discover that we are a term in a relation.
Ancient scepticism discovered it. And we shall bring forward some examples
of how in more recent times one has more and more come to see the
significance of this relation for all our knowledge. What particular part
the subjective relation (as we may call it for brevity) plays in the
individual case cannot of course be said in advance. It is through analysis
of the advancing knowledge of the world that the share which man (the
subject) has in this knowledge can gradually become clear. It is some
examples from the history of the more recent age's natural science and
epistemology that shall now be brought forward.

\subsubsection*{β.~On Copernicanism.\footnote{In this and the next section I
build on my presentation of the history of the more recent philosophy.}}%
\addcontentsline{toc}{subsubsection}{β.~On Copernicanism}%
\label{sec:3d-beta}%
% pp. 62--67

\noindent One of the grounds on which Copernicus built his hypothesis was
the subjective concept of relation. Sensation cannot, he says, immediately
tell us what it is that moves --- whether it is the sensed thing, or the
sensing subject, or perhaps both with different velocity or in different
direction. If now the earth moved, the sensing subjects would move too, and
there would be no ground for ascribing motion to the sun. A century before
Copernicus a similar train of thought had already been asserted by the
profound thinker Nicolaus Cusanus, who drew attention to the fact that
wherever a man is thought to be, he will believe himself to be at the
centre, no matter whether he is somewhere on the earth or on the sun or on
other stars. One has therefore no right to speak of a midpoint in the world,
and thus no right either to assert that the earth is this midpoint. And from
this it again follows that one has no right to assert that the earth is at
rest. When we notice no motion, this may after all come from our not having
any absolutely resting point in relation to which it could come forward.
This train of thought is in Cusanus a member of a whole way of seeing which
impresses the determination through relation (relativity) of all thoughts,
and this way of seeing is again grounded in the fact that all knowledge comes
about by the connexion of a plurality of objects, which connexion
necessarily brings these objects into relation to one another and to
reciprocal determination by this relation. Cusanus has on the whole a
remarkable capacity for thinking every object in its definite relations. He
thereby stands as a precursor of Copernicus, who showed what new world-picture
could be the consequence of this way of seeing, and of Giordano Bruno, who
like Cusanus derives the concept of relation from the nature of our knowledge
(del modo nostro de intendere). Our thought constantly works by determining
one object by another, by setting limits and cancelling limits --- in virtue
of one and the same principle. Thus the one place is determined in relation
to the other, and motion in relation to a presupposed fixed point. If I seek
such a point on the sun, a motion will look otherwise than when I seek it on
the earth. The old world-picture presupposed what was precisely to be
proved, that the earth was the one fixed point in relation to which all
motion was to be determined in order to be true. Of centre, pole, zenith and
nadir we can speak absolutely only when we assume the earth to be at rest.
Of particular interest is that for Bruno the relativity of time follows with
that of place and motion. Aristotle had (Physica IV, 14) taught that the
concepts time and motion hang closely together and are reciprocally measured
by one another, and for him the measure was given in the regular motions of
the heavenly bodies. To this Bruno remarks\footnote{\textit{Acrotismus},
Art. 38 (Opera latina. Ed. Fiorentino. I, p.~143---146).} that since one and
the same motion must look different from the different globes in space,
there must be as many different times as there are stars! --- Interesting in
Bruno is on the whole the way in which the concept of relation leads him to
demand new observations and new experiments where one had hitherto rested
content with the immediately given as containing its explanation in itself.

In Galileo there enters the turn that a sharp distinction is drawn between
the mathematical and the sensuous apprehension. In mathematical knowledge
the difference between divine and human thought falls away for Galileo. It
is his studies of Plato that are at work when he assumes that in mathematics
the difference between existence itself and our apprehension does not hold:
here human knowledge shares in the necessity with which the Deity thinks the
truths that bear the connexion of existence; only this difference remains,
that what men must laboriously and successively work to understand comes
forward for the Deity in immediate intuition, and that the compass of divine
knowledge far surpasses that of human knowledge. Only for the mathematical
(divine) apprehension is there absolute space and absolute motion. For our
sensuous apprehension all motion is relative, since after the overthrow of
the ancient, bounded world-picture we have nothing absolutely unmoved in
relation to which motion can be established. Our sensuous apprehension is
conditioned by our standpoint on the earth: ``Think the earth away, and
there is no sunrise or sunset, no meridian, no day and no night!'' And here
on earth it is again our bodily organism that conditions the properties we
--- apart from the purely mathematical ones (primi e reali accidenti) ---
ascribe to things, thus all the proper sense-qualities (colour, smell, taste
and so on). Take the body away, and all sense-qualities fall away too, and
then our knowledge will be able to coincide entirely with divine knowledge.

The sharp opposition between mathematical (divine) knowledge, for which
there is absolute space and absolute motion, and sensuous (human) knowledge
is the more remarkable to find in Galileo, since his great deed was
precisely the methodical co-operation of deduction and induction, of
mathematics and experiment. But this opposition gained great influence in
the following centuries, especially as it was supported by Newton's great
authority.

On account of the proposition of the relativity of motion, the celestial
phenomena could purely formally (mathematically) be explained just as well
by the sun going round the earth as by the earth going round the sun. And
Tycho Brahe had set up an intermediate form between the two hypotheses,
according to which the earth stood fixed while the sun and the planets went
round it. Gassendi instituted a comparison between the three hypotheses thus
set up and sought to show that purely scientifically one could not choose
between them. But then Newton showed in his famous work that the laws Kepler
had discovered for planetary motion could be derived if one extended the law
of gravitation to hold for the heavenly bodies, and then it must necessarily
be the earth and the planets that moved about the sun. In virtue of the law
of falling motion, then, only one hypothesis was at last possible. In an
analogous way Leibniz at about the same time led beyond the scepticism which
the relativity of motion apparently motivated, in that he maintained that
what was decisive for knowledge was not motion itself but the laws of
motion, and here he then set up --- in criticism of Descartes' doctrine of
the persistence of motion --- the proposition of the persistence of force.
Here then was a reality that did not conform to the observer's standpoint in
space. The two great investigators led, each by his own path, beyond the
apparent subjectivism and took the decisive standpoint in the laws of
motion.

Nevertheless Newton maintained (here surely under the influence of the
Platonist Henry More), like Galileo, a sharp opposition between sensuous
space, to which the popular apprehension holds, but which makes no absolute
determination of place possible and exhibits only relative motion, and
absolute space, which stands in no external relation whatever, and which
permits absolute determination of place and absolute motion. It is in
scientific orientation in the world (in rebus philosophicis) that we make
use of the absolute places (loca primaria), and it is then mathematical
space that lies at the basis. In practical life (in rebus humanis), on the
other hand, we content ourselves with sensuous space, which makes no
absolute determination possible. This conception, which hangs together with
Newton's peculiar metaphysics, became for the time being the dominant one
within natural science. It was a formidable barrier that was by the great
naturalist's doctrine on this point set against the continued application of
the concept of relation. Since one no longer, as Copernicus still did,
assumed the sphere of the fixed stars (``the eighth sphere'') to stand
fixed, so that it --- as communis universorum locus (to use Copernicus'
expression) --- could serve as basis (point de repère) for absolute
determination of place, it was really a purely mystical reference to appeal
to absolute space, which was not accessible to any human experience. --- At
this point of the discussion two philosophical thinkers now intervene in a
peculiar way.

\subsubsection*{γ.~Philosophical Standpoints.}%
\addcontentsline{toc}{subsubsection}{γ.~Philosophical Standpoints}%
\label{sec:3d-gamma}%
% pp. 67--71

\noindent It was a new Platonism that triumphed with Galileo and Newton, in
peculiar opposition to the way in which these great investigators in their
natural science applied the reciprocal action between thinking and
observation. But an opposition now arises which may be compared with that
raised in antiquity by Carneades and Aenesidemus against the Platonic
doctrine of Ideas. It was Berkeley and Kant who raised this opposition. They
both criticize, each in his own way, absolute space, and maintain that if
one leaves the knowing subject out of consideration, all truths and all
riddles fall away at the same time.

Berkeley's principal thought was that we think and must think in examples.
General concepts, propositions and laws we know only from the special cases
in which they display themselves, and it is unjustifiable to ascribe to them
and to the relations they express existence as a world of their own behind
the world of qualities and differences of quality appearing to observation.
The ``primary and real'' properties, to which Galileo ascribed absolute
existence, we nonetheless always apprehend only under definite conditions:
at some distance or other, macroscopically or microscopically, and so on.
Extension in space must share the fate of all other properties. Every motion
must have a definite direction and a definite velocity. There is neither
extension nor motion in complete generality. When the question is raised
about the right place, the right magnitude, direction or velocity, such a
question can be decided only with regard to the apprehending and thinking
subject. If I, says Berkeley, being the childlike believer he was, had been
present at the creation, I should have seen things come into being in the
order the Bible relates. But in daily speech we disregard the subject,
although it must always be understood. --- Here the opposition to Newton
appears clearly; Newton thought that in strict science we disregard the
subject. Berkeley maintains that in science (in rebus philosophicis, to use
Newton's expression) a subject must be understood just as much as in
practical life (in rebus humanis).

Kant has, as already mentioned, seen that relation is a fundamental concept
holding for all knowledge, only that he, with decided inconsistency,
reserved a relationless ``Ding an sich'' to explain that there exists for us
something we have not ourselves brought forth.\footnote{When ``Ding an
sich'' is to contain the ground of matter (and form) in our knowledge, it
nonetheless stands in relation to knowledge, notwithstanding that it was to
be relationless. Or our experience stands related to it, but it not to our
experience. It is the same one-sided relation that Thomas Aquinas tried to
carry through in theology. See above p.~47 Note.} Among the relations it is
for him, as for Berkeley, the relation subject---object that occupies the
interest, and he maintains the fundamental position of the subject in this
relation. Even in pure mathematics and exact empirical science a knowing
subject is presupposed. Only for such a subject is there experience in the
strict sense of the word, an experience in which mathematical concepts find
their application and have their real significance.

Kant began as an eager Newtonian and for the time being maintained, like his
master, absolute space as an objective presupposition for strict empirical
science as well as for pure geometry (which for him had attained its
definitive foundation with Euclid). Even in ``Kritik der reinen Vernunft''
he maintains an absolute space, only that the Newtonian space has so to
speak been ``turned inward'' and become a human form of intuition instead of
a divine form of intuition, a sensorium hominis instead of, as for Newton, a
sensorium dei.\footnote{Cf.\ \textit{Om Kontinuiteten i Kants filosofiske
Udvikling}, p.~31.} Of particular interest, however, is that he emphasizes
that the concept of space (and the same holds of the concept of time) is not
merely a general concept, but what one might call a typical individual
concept: every definite space is a part of space, not merely an example of
the concept space. And it must also be noted that when he calls space a form
of intuition, he does not mean by intuition a result finished once for all,
but a function, an intuiting, which leads to the forming of intuitive images
or the construction of schemata. On the whole Kant means by ``forms'' ways
in which the connecting function, the synthesis in which all consciousness
and all knowledge consists, takes place in the individual cases. We
construct space by adding point to point, piece to piece according to a
definite rule. It is in the identity of such a rule that the unity of
consciousness expresses itself in the individual case.

By a closer carrying through of this train of thought Kant later came to
deny the subjective reality of absolute space, just as he had already in
``Kritik der reinen Vernunft'' denied its reality as existing in itself.
Already in the first edition of his principal work he had stated that, apart
from all things that fill and limit, absolute space is nothing other than
``die blosse Möglichkeit äusserer Erscheinungen''. Some years later (in
``Metaphysische Anfangsgründe der Naturwissenschaft'') he goes a step
further, relying on the relativity of every place and every motion. In
absolute space it should indeed be possible to disregard all relations. But
when do we have done with setting all points in relation to one another,
with drawing all lines and so on? Where is the point or the limit which,
itself relationless, determines all points and limits? Kant therefore
declares that pure, absolute space ``für alle mögliche Erfahrung nichts
ist'': ``Der absolute Raum ist an sich nichts und gar kein Objekt, sondern
bedeutet nur einen jeden relativen Raum, den ich mir ausser dem gegebenen
jederzeit denken kann.'' (Anf. Skr. 1786, p.~16 cf.\ p.~3). Kant here in
fact espouses Berkeley's doctrine that we think in examples. Absolute space
means merely the negative fact that no given space is conclusive. New spaces
can always be imagined. It is the very thought of relation, of
relation-setting, that leads Kant away from his absolute form of intuition.
Space becomes an ``Idea'' in Kant's sense, a boundary concept which rejects
conclusion and maintains the possibility of continued formation of series.
The concept of space acquires only regulative significance, like the
``Ideas'' set up by Kant himself, and at bottom it belongs under ``the
cosmological Idea'' (the concept world).

Kant has himself compared his philosophy with Copernicanism. And in a way he
and Berkeley might have been discussed in the section that dealt with the
thoughts developing ``on Copernicanism''. But the two philosophers were
nonetheless mostly occupied with combating the dogmatism which, with
supposed support in mathematical science, had maintained itself in spite of
the victory of Copernicanism. And as regards Kant in particular, his
principal endeavour was at one and the same time to maintain the great
significance of the Newtonian concept of science and to lead the absolute
forms Newton had still let stand back to being forms of the working thought.

Although Kant constantly set time parallel with space, he does not for time
draw a similar consequence as with respect to space. And yet it would ---
particularly with his remarkable doctrine of the ``schematisms'', which
emphasizes the activity in the intuition of time --- have lain near him to
see that we have as little done with time as with space, and that absolute
time can mean only that new relations of time are always possible. ---

I need not here enter upon a criticism of Kant. Such a criticism would
particularly aim at showing that the very concept of experience in Kant's
strict sense indicated above is an ``Idea'', like space and time, and that
Kant can assume a ``Ding an sich'' without relations only by a strange
inconsistency.

\subsubsection*{δ.~The New Theory of Relativity}
\addcontentsline{toc}{subsubsection}{δ.~The New Theory of Relativity}
\label{sec:3d-delta}
% pp. 72--77
%
% REDONE 2026-08-23 from the completed transcription. The earlier version of
% this section was made from a pre-scan Danish text and diverged from the
% printing in eleven places; all are corrected here. In particular:
%   * Six runs of \emph{} that the printing does not have are now plain roman:
%     the Neumann and Mach titles, "points de repere", "ultimate character",
%     "plus ultra", "Tidsskrift for Fysik".
%   * Small caps restored at first mention, as the printing sets them and as
%     the companion volume's translation does throughout: Neumann, Mach,
%     Painleve, Einstein, Langevin, Whitehead, Haldane, Nicholson, Kroman,
%     Lucien Fabre.
%   * The Mach citation takes its full stop BEFORE the parenthesis, as printed.
%   * The Einstein note keeps its comma inside the italic and ends without a
%     full stop, as printed.
%   * Three footnote anchors on p. 75 stand before the comma/period, as printed.
%   * The section runs to p. 77, not p. 76: the Campanella paragraph spans the
%     page turn and the section 4 head does not open until partway down p. 77.

\noindent Despite the great expansion of the horizon that Copernicanism
gave impetus to, natural science had held fast to the concepts of absolute
place and absolute motion, which had their justification only within the
old worldview with its bounded and fixed frame. Galileo's and Newton's
absolute space replaced the space bounded by the eighth sphere. But there
came a time when the concept of relation here too made its consequences
felt, and not only with respect to the relations among objects, but also
with respect to the relation between objects and the knowing subject.

It was the first proposition of kinematics, and thereby of natural
science, that was first examined from the standpoint of the concept of
relation. \textsc{Neumann} (Ueber die Prinzipien der Galilei-Newton'schen
Theorie, 1870) drew attention to the fact that the law of inertia,
according to which the velocity and direction of a body change only through
external influence, necessarily presupposes in each individual case the
specification of the particular place in relation to which velocity and
direction are preserved when no external influence is exerted. Neumann
could find no place that could lay claim to being taken as the starting
point for an absolute coordinate system, but held that one must
hypothetically assume a center in order for the concept of absolute motion
to retain its meaning. \textsc{Mach}, who at roughly the same time had arrived by
his own path at the insight that the law of inertia rested on the
presupposition of a determinate coordinate system, did not find Neumann's
solution satisfying. He demanded that the concept of relation be carried
through. The dogmatic conception of the law of inertia, in his view, only
falls away when one becomes clear that in itself any body in the world can
be used as a starting point, and that in practice in each individual case
we use one body or another as the presupposed fixed starting point for the
coordinates in relation to which the velocity and direction of a motion are
determined. (Die Erhaltung der Arbeit, 1872). This is the path that
has in fact always been taken, only within the old worldview, which rested
precisely on the presupposition of an absolute center and an absolutely
fixed boundary-sphere of the world, no choice was required. Mach lays
great weight throughout on the historical presuppositions on which the
research of each period rests, but then shows how these presuppositions
vary from age to age, and that none of them can lay claim to absolute
necessity and intelligibility. What matters here always comes down to what
French physicists call points de repère.\footnote{Cf.\ e.g.\
  \textsc{Painlevé} (\textit{Mécanique}, p.~399) [in the collective work
  ``De la Méthode dans les Sciences'', 1909]: ``\ldots\ les mouvements
  absolus, c'est-à-dire au fond les mouvements convenablement repérés.''
  [i.e., absolute motions are, at bottom, motions suitably referred to
  reference points.]}

An even more radical application of the concept of relation appeared at
the beginning of the twentieth century. What is of interest in this
connection is that it has proved necessary to take special account of the
apprehending subject and its standpoint. Copernicus had already pressed
this point as regards the apprehension of the motions of celestial bodies,
and Neumann and Mach as regards the velocity and direction of motion.
\textsc{Einstein}\footnote{\textit{Über die spezielle und die allgemeine
  Relativitätstheorie,} (1916)} has now sought to show that what holds
for space and motion holds also for time. That this was not discovered
earlier is, according to Einstein, to be explained by the fact that
attention was paid only to velocities that were very small in relation to
the velocity of light. The final temporal measure becomes, for modern
physics, the velocity of light, whereas for Aristotle, and even for the
Copernicans, it was drawn from the motions of celestial bodies. (Bruno's
doctrine of the relativity of the concept of time rests on this measure.
See above.) On the basis of the electromagnetic theory of physical
phenomena, for which the mechanical conception of nature stands as only a
special case, Einstein seeks to show that there can be demonstrated
circumstances under which two observers, of whom one is attached to fixed
and the other to moving axes, but both at rest relative to their own axes,
will not arrive at the same apprehension of time: what for one observer
stands as a relation of simultaneity between two events, will for the other
stand as a relation of succession. A temporal specification acquires
meaning only when it is stated in relation to which coordinate system it
holds, and in particular whether that system is itself relatively at rest
or in motion.

If one were to hold that an illusion must be present in one or the other
observer, the answer from the Einsteinian standpoint is that each observer
has his own system of measurement, which is just as legitimate as
another's, though it is of course perfectly possible for one observer to
attempt to place himself at another's standpoint and reason from
it.\footnote{\textsc{Langevin} in ``Bulletin de la Société française de
  Philosophie'' 11 October 1911. --- It must presumably be assumed that
  the observer who can place himself at another observer's standpoint knows
  sufficient physics.}

The facts and inferences on which the Einsteinian theory rests are still
a matter of dispute, and views stand sharply opposed to one another. In
such a dispute the philosopher as such cannot intervene, and philosophy
has time to wait. Whether physics can really demonstrate situations in
which what stands as simultaneity for one observer stands for another
observer as successive, physicists must determine among themselves. If
this must be answered with yes, we stand before the significant result
that natural science has once again, as with the Copernican hypothesis,
by its own path come to see the necessity of not merely holding to the
content of observations, but of taking account of the observer's
presuppositions. In an interesting discussion in the ``Aristotelian
Society'' (1919), \textsc{Whitehead} impressed upon us the significance of
what he aptly called ``the percipient event''\footnote{See also his
  \textit{Principles of Natural Knowledge} (1919), p.~68. --- Whitehead's
  conception differs from Einstein's, however, in that ``the percipient
  event'' for him is only one factor alongside many others, and not, as in
  Einstein, of essential significance for all measurement of time and
  space. Cf.\ on this point \textsc{Haldane}: \textit{The Reign of
  Relativity} (1921), p.~69\,f.; 107\,ff.},
% The Danish note here prints "Belydning" for "Betydning" — a broken type,
% logged as a printer's defect at its site in transcription.tex. The error is
% in the Danish glyph and cannot be carried into English, so the intended
% sense ("significance") is translated and the defect noted here instead.
% The footnote anchor stands BEFORE the comma, as printed on p. 75; the same
% holds at the two other anchors in this paragraph.
and \textsc{Nicholson} declared that among the data that are of decisive
significance (ultimate character) there must also be counted ``the
essential framework of our modes of perception.''

Against the Einsteinian hypothesis, as against earlier demonstrations of
the epistemological significance of the concept of
relation,\footnote{Cf.\ \textit{Den menneskelige Tanke} [The Human Mind],
  p.~304.} it has been urged that there can only be one truth in each
individual case, and that it would contradict the principle of
non-contradiction if two observers, of whom one declares simultaneous what
the other declares successive, were both to be right\footnote{\textsc{Kroman}
  in Tidsskrift for Fysik, XIV, p.~15.}. But there is no reason to fear
for the unity of truth, even if Einstein should be right. The real truth
lies in this case in the necessity, physically demonstrated, that different
observers under certain determinate conditions arrive at different
apprehensions of time. Each individual observer apprehends time as he must
under his conditions. But if an observer is also a physicist, he will
(assuming the correctness of Einstein's theory) be able to explain to us
why another observer arrives at a different apprehension than he himself
does. (Cf.\ Newton's and Leibniz's emphasis on the lawfulness of motion as
the expression of its objectivity. Above, p.~66.) Naturally, the physical
explanation of the different observers' different apprehensions of time
itself rests on certain presuppositions, presupposes a coordinate system
from which the coordinate systems of the various observers are determined,
and here there is the possibility of a new discussion. When this discussion
comes to an end, it is because no standpoints can be demonstrated that are
still more fundamental. Kant long since maintained that all necessity is
hypothetical. But it is necessity nonetheless --- or rather: precisely
because of this it is necessity. A necessity without presuppositions we do
not know.

If the new theory of relativity wins continued confirmation, what comes to
hand will not merely be a result of physical research that is interesting
in and of itself. This theory also contains, as Einstein himself has
declared, a spur to continued inquiry. If it is the case that there is the
possibility of different systems of relations in the apprehension of an
object, this constitutes precisely an invitation to discover all the
particular relations that lie at the basis in each individual case, and
inquiry is thus led ever further rather than closing off in dogmatic
fashion, as if the presuppositions one had hitherto worked with were the
only possible ones. There are, from an epistemological point of view, two
interests that attach to the concept of relation. One is that light is
cast on the nature and conditions of our knowledge; the other is that new
tasks are indicated, a plus ultra appears, in that every thought
points beyond itself. ---

Campanella said in his apology for Galileo that if Galileo were right,
one would have to philosophize in a new way. But if Einstein is right,
one need not therefore philosophize in a new
way.\footnote{\textsc{Einstein} himself agrees with this. In his preface to
  \textsc{Lucien Fabre}'s book ``Les théories d'Einstein'' (1921) he says: ``The theory
  of relativity neither can nor will provide a world-system, but only
  indicate a limitation to which the laws of nature are subject.'' That it
  is not merely a matter of a limit, we have seen above.}
Even from a purely philosophical point of view, as the foregoing
presentation has shown, attention was long since drawn to the significance
of the concept of relation for knowledge and worldview. But if Einstein is
right, there now lies before us a new, unexpected, but significant
experience in that respect. ---

Through the necessity of taking account of the knowing subject, and
through the possibility of a conflict between the each-for-itself
necessary apprehensions of different knowing subjects, the newer natural
science --- in Einstein as in Copernicus --- offers an analogy to what
asserts itself powerfully in the domain of the human sciences, as I shall
now attempt to show.


% ============================================================
%  4. THE CONCEPT OF RELATION IN ETHICS
% ============================================================

\section{The Concept of Relation in Ethics}
\label{sec:4}
% pp. 77--98

\subsection{Historical Research and the Concept of Relation.}
\label{sec:4a}
% pp. 77--81
% Section 4's lettered heads are letterspaced and RUN IN (the text follows on
% the same line after an em dash), unlike section 3's, which stand on their own
% line. The transcription drops the run-in dash and sets them as headings;
% that treatment is followed here.

\noindent Aristotle taught, as already mentioned, that only the general, not
the individual as such, could be an object of science. This hangs together
with the logical-schematic character of his philosophy. Science was for him,
as for Plato, subordination of the particular and individual under general
concepts. But now for Aristotle everything real was at the same time
particular and individual; it is on this point that he comes into opposition
to Plato. The difficulty arising hereby, and which has been brilliantly
illuminated by \textsc{Zeller} (Die Philosophie der Griechen, II, 2), can be
overcome only through a closer discussion of what one understands, and
scientifically must understand, by reality; and such a discussion was, as we
have seen, introduced with great energy by the Academics and Sceptics of
antiquity.\footnote{It is another matter how the difficulty could be cleared
up within the Aristotelian philosophy itself. \textsc{Hans v. Arnim}
(\textit{Die europäische Philosophie des Altertums}. Kultur der Gegenw. I,
6, p.~173 f.) finds the explanation in the fact that both species and
individuals are in Aristotle regarded as ``beings''. But thereby a new
question is merely called forth, namely how and why the one kind of
``beings'' can be an object of science and the other not.}

Already here the thought came forward that the criterion of reality can
consist only in the fixed connexion between observations. Just as the formal
sciences spring from a striving to preserve agreement with oneself, so the
real sciences spring from a striving to find agreement between the objects
given to observation. And the criterion of reality lying in this must hold
just as well for the individual, that which perhaps occurs only a single
time, as for that which can repeat itself again and again because it holds
for whole series of individual objects.

An object is known as real, partly through its connexion with other objects
concerning which no doubt has been awakened, partly through the connexion
that can be found between the various forms under which the object itself
appears at different times, and between the various properties with which it
appears at a given time.

Historical research now aims precisely at both kinds of connexion. The
events, personalities or institutions that are its objects it investigates
in relation to the circumstances under which they are present, and each
single one of those objects it investigates with regard to its inner
connexion. In both respects a beginning is made with criticism of reports,
but this criticism itself lays the named criterion of reality at the basis.
Thereafter it becomes the task to find the bond connecting the object with
other objects, and the bond that ties together the object's various sides
and stages.

The object of science is neither the general nor the individual in and for
itself alone, but reality, which can be known only through the bonds uniting
objects among themselves as well as the elements within the single
object.\footnote{Cf.\ \textit{Totalitet som Kategori}, p.~59---66.} What
cannot be illuminated through application of the concept of relation cannot
be understood.

Philosophers of history have in the most recent time brought forward three
concepts as particularly characteristic of historical inquiry\footnote{Anf.
Skr. p.~71---72.}, and these three concepts express precisely relations
between an object, its conditions or its effects.

Historical centre is the definite event, personality or institution whose
origin, subsistence and development the historian wishes to investigate.
About it one gathers everything that can contribute to its illumination, but
takes no account of what neither furthers nor hinders it. Correlative to
this concept is the concept historical threshold, which denotes the limit of
what can be included in the individual investigation. Historical greatness
denotes that character of an event, a personality or an institution whereby
there are expressed in them forces working in the direction of the solution
of new, significant social tasks.

Now there stands in history as in nature a struggle for existence. The
historian stands before events and personalities in the human world as the
biologist stands before plants and animals in their struggle for existence.
The point in both cases is to understand the contending parties each for
itself, as regards their nature and their conditions, and thereby to
understand how a struggle arises between them. The biologist as such does
not ``side with'' any of the organic forms, even if his intellectual
interest has made a single one of these forms the centre of his
investigations, or even if he can feel a certain sympathy for certain forms
of life, as when it gladdened Darwin when he thought he could find functions
in plants that might point to a higher stage of life than one was accustomed
to ascribe to them. For the biologist the point is merely to understand the
necessity of the struggle under the given conditions. Nor does the historian
as such take sides, but seeks to find what conditions assert themselves in
each of the contending parties, in order to understand how interests,
thoughts and feelings have had to unfold themselves in a definite degree and
in a definite direction in each of them, until the collision became
unavoidable. Every nation, says William James in one of his letters, has its
ideals, which are a dead secret to other nations, and must under the
influence of these ideals develop in its own peculiar way. If this is so,
there is required of historians an understanding both of how these ideals
themselves have developed and of how they have determined the national
development. Thereby it again becomes possible to understand the tension and
the strife between nations. It has been said of Leopold Ranke that when in
his historical inquiry he avoided one-sided partisanship, it was not because
he was neutral, in the sense of indifference towards all contending parties,
but because he was universal and could enter into and follow the life
stirring in each of them. This historical universality is akin to the
artistic universality that enabled Shakespeare to let each of his figures
unfold according to their own inner laws. His impartiality towards persons
and fates sprang not from indifference and remoteness, but from being seized
by life as it can stir in different men and lead them to collisions that
perhaps become tragic.

The historian (and the poet) stands here in principle as the physicist
stands when he establishes certain relations under which what for one
observer stands as simultaneous stands for another observer as successive.
% The Danish here prints "hvike" for "hvilke", the l dropped; logged as a
% printer's error at its site in transcription.tex. The slip is in the Danish
% orthography and cannot be carried into English, so the sense is translated
% and the defect noted here.

\subsection{Ethics and the Psychological Relations.}
\label{sec:4b}
% pp. 81--87

\noindent Historical valuation is not properly an ethical valuation. It
concerns the significance of an event, a personality or an institution for
the continued actual development, but does not pass its judgments from an
ethical ideal. With great fineness and energy \textsc{Heinrich Rickert} has
in his work ``Grenzen der naturwissenschaftlichen Begriffsbildung''
demonstrated this difference between history and ethics. And from the
historical side it has been energetically stated that historical inquiry
does not penetrate to the place where actions and endeavours have their
source: ``Das tiefinnerste Geheimnis einer Seele zu finden, damit deren
sittlichen Wert, das will sagen, den ganzen Wert der Person richtend zu
bestimmen, hat die historische Forschung keine Methode und keine
Kompetenz''. (\textsc{Droysen}: Geschichte Alexanders der Grossen, p.~259).

If nevertheless the ethical standpoint of value involuntarily forces itself
forward in the historian when he tries to give his presentation its final
rounding-off, he will stand before the difficulties of principle presenting
themselves in every attempt to ground ethical judgments, thus in every
ethics.

At this point a new class of categories appears, which does not coincide
with any of the classes --- the fundamental, the formal, the real --- that
we have hitherto discussed: the concepts of value or the ideal categories.
The concept of relation asserts itself here at once, in that value can be
determined only in relation to value. All willing is a preferring, and every
feeling has its degree and its quality determined by its relation to other
feelings; therefore in every valuation everything comes down to what stands
beforehand or simultaneously as valuable. At last all valuation in the
individual cases will rest on the relation to a fundamental value. Just as a
place can be determined only by its relation to other places, a time only
through its relation to other times, a reality only through its relation to
other realities, so every determination of value presupposes that certain
other values stand fixed beforehand. The fundamental value corresponds to
the globe or the coordinate system from which we orient ourselves in space,
--- to the event in relation to which we date other events, --- to the
reality which in the individual case decides what is to be reality for us.
And if one investigates more closely the various possible fundamental values
asserting themselves in human valuations, one will find that they all
presuppose some whole or other whose subsistence and development comes to
stand as end and thereby determines the valuation. Such a whole may be the
valuing personality itself, or a society (family, estate, state, church), or
a sum of endeavours (art, science and so on). A valuation will at last be
determined by an action's relation to the conditions of such a whole's
subsistence and development. Only when a whole given in experience is
presupposed, whose conditions of existence and development can be discussed,
is a rational valuation, and thus a rational ethics, possible.\footnote{On
the connexion between the concept of totality and the problem of valuation
see: \textit{Den menneskelige Tanke}, p.~237---241; 346---369.
\textit{Totalitet som Kategori}. Kap.~5.} Ethical idealism often works
against itself by seeking beyond all relations. In the several places where
Plato discusses ``the Idea of the Good'', it is expressly raised (like other
Ideas) above the world of relations. He does indeed maintain (especially in
``Gorgias'' and ``Philebus'') a connexion between the valuable and definite
relations of measure, namely proportionality (``geometrical likeness'') or
harmony between fullness and limitation. But such relations of measure hover
in the air when they are not derived from the conditions of life of a
certain definite whole, and here Plato oscillates between an individual and
a social totality.

The concept norm concerns the means by which a presupposed whole can subsist
and develop. The presupposed whole then stands more or less consciously as
end, in that its value denotes a fundamental value. The concepts value, end
and norm lie in the same series, and in a rational ethics they must lie in
this definite order. No ethics can withdraw itself from the connexion
between these three concepts. Often one conceals the first term in the
series and holds either to the concept end or to the concept norm, as though
they were primary. We meet here again that lack of consciousness of where
the one compass-leg stands which so often leads to a misjudging of the
significance of the concept of relation. One then regards the position of
the other compass-leg as absolutely given, self-evident or ``objective''.
Not least does this appear in moralists who think they have secured
themselves an ``objective'' standpoint by proceeding from ``society''
(without even saying what society they mean). Or it is, as in Kant, the
concept norm that one lays at the basis. Kant, who in his epistemology is
clear that every necessity is hypothetical, maintains an ethical necessity
(a categorical imperative) which is not again to rest on any presupposition,
and in various forms later moralists have followed him in this. When an
ethical system is criticized, the criticism will --- apart from an
investigation of the system's formal consistency --- essentially aim at
showing that the system presupposes a fundamental value which is not itself
drawn into the investigations. Criticism will seek the hidden compass-leg.
Even he who proclaims a ``revaluation of all values'' must presuppose some
value or other in relation to which he undertakes the revaluation.

The fundamental value will at last be determined by the impulse to
self-preservation in the individual (in so far as it can be thought
isolated) or in the greater or smaller society with which the individual
feels himself in solidarity, so that this society's weal and woe, and
thereby its ends, become the individual's too and determine his ethical
tasks. There may, however, also be expressed a striving to make the single
moment or a single side of the individual's life a whole in and for itself.
From the single moment to the human race in present and future there is a
series of possible foundations for ethical systems, and these foundations
are, as the history of ethical ideas shows, in part made use of.

The comparative moralist now has (as in their domains the physicist and the
historian have) the task of understanding the conceptions which the various
foundations (fundamental values) make possible, and which may come into
decisive conflict with one another. Thereby is understood the fact that the
same action, the same personality, the same institution can be the object of
conflicting ethical judgments. On account of the lawful connexion which can
be demonstrated by psychological and historical means in the spiritual
world, we seek an explanation of why others judge otherwise than we do, and
thus have other values, ends and norms than we.

And it is not merely the comparative moralist who has use for such
understanding. For it is after all a concern of every ethical standpoint to
get men to know and acknowledge the values, ends and norms which it itself
lays at the basis. And a beginning must hereby be made by taking men where
they actually stand, entering into their presuppositions in order to be able
to lead them to new presuppositions. Socratic irony is, seen from an
essential side\footnote{Cf.\ \textit{Den store Humor}, p.~63.}, a means of
getting men to open themselves, to reveal themselves, in order that it may
thereby become clear how help, education and co-operation can become
possible. It is the analytical method which alone can lead from one
standpoint to another, when this is possible at all. It is of course easier
to begin with the construction of the ``right'' ethics and then leave it to
men how they can come to terms with it. But every value and every valuation
has in virtue of the concept of relation its definite psychological
presuppositions, and the point is first and foremost to discover how far
these are present, so that a common valuation can become possible. Only a
common fundamental value makes common valuation possible.

Misjudgment of the relation between what is a moralist's fundamental value
and what he is from this foundation led to teach has often led to unjust
judgments about ethical thinkers. Thus it has been thought that the weight
which philosophers like Adam Smith, Helvetius and Bentham have laid on the
acquisitive impulse, on individual interest and on the impulse to happiness,
must be due exclusively to a preoccupation with individual or egoistic
interests, so that their own fundamental value must consist in such
interest. A closer investigation shows, on the contrary (as I have shown in
the relevant section of ``Den nyere Filosofis Historie''), that the named
thinkers all proceed from universal sympathy and from conscious regard for
the welfare of human societies. Their conception is that society will thrive
best by each individual building on his experiences of what serves him best,
and that no authority standing outside him can know this better than he
himself, when his eyes are opened. Through free labour and through the
reciprocal action of independently striving individuals, the fundamental
value on which they build will for the named thinkers best come into its
own.

The individual moralist must be conscious that the fundamental value which
bears his system can itself be brought forward for testing. If the system
has a vulnerable point, it will at last lie here, even if the system is
otherwise developed with ever so great consistency and supported by ever so
many experiences. And he must, when he is sufficiently critical, allow Henry
Sidgwick to be right that it --- at any rate at certain stages of
development --- is good that in a society different, even conflicting,
ethical conceptions should be asserted, even though he can of course follow
only his own.\footnote{\textit{Methods of Ethics,}\textsuperscript{2}
p.~450.} In the ethical domain too, truth can work its way forward only
through strife. Psychology and history have then to see to it that behind
the burning strife there can be an underground understanding.

The whole conception of the position of ethics which is developed here and
in what follows in connexion with the concept of relation, I have maintained
from 1887 onwards in my various writings, and it has constantly won new
firmness with me, not least on account of the criticism to which it has been
subjected from various quarters.

\subsection{Ethics and the Historical Relations.}
\label{sec:4c}
% pp. 87--90

\noindent The working out and carrying through of ethical norms --- thus the
determination of the subordinate ends which the principal end leads one to
set, and in which the derived values that the fundamental value gives rise
to find their expression --- is possible only in reciprocal action with the
given historical conditions. Neither the valuable qualities of character
(``virtues'') nor the social arrangements which a fundamental value requires
for its full realization can be constructed in advance; they cannot be
``introduced'', as in old days one thought one could ``introduce'' a
religion or a constitution. In any case the historical conditions will at
different times and in different places give both individual qualities of
character and social arrangements essential peculiarities which cannot be
derived solely and alone from the general norms to which the fundamental
value gives rise. Even a brief survey of the history of the virtues, such as
the one I have attempted to give in my Ethics (X, 4), shows this. As regards
the social arrangements, it is seen from the history of marriage (Etik XV)
and of the social question (Etik XXV). The old leather bottles cannot at
once be replaced with new ones. Besides, ethical norms always have a general
character, and the specialization can be grounded only on the historical
conditions under which they are to be carried through.

Abstractions and vague dreams are not enough. The point is for the time
being that the direction for the development of character and the ordering
of society should be determined by the fundamental value. A complete
programme only experience can lead one to work out. It goes with ethical
norms as it goes with the concept of cause in epistemology, which likewise
does not in itself show any solution of special questions, but merely
motivates an asking and a seeking which acquires a provisional character at
various stages of science's historical development.

There lies here the possibility of tragic conflicts, in that those who are
animated by a fundamental value may lack an eye for the historical
conditions, while those who stand firmly on the ground of historical reality
may lack an eye for values that demand realization. In the present day there
stands more than ever a struggle between social ideals and social realities.
Bertrand Russell, for whom the social ideal is given in communism, finds
that the price which by the application of Bolshevik methods is to be paid
by the human race in order to reach communism is far too high and will even
cancel the value which the carrying through of communism would in itself
have. A principal cause of this he finds in the fact that in Bolshevism
hatred of the old plays a greater part than hope of the new; these
psychological moments too are historically conditioned.\footnote{\textsc{Bertrand
Russell}: \textit{The Practice and Theory of Bolschevism}. (1920),
p.~146---178.} The counterpart to this is found more than sufficiently
represented in Europe as hatred of the new and hope of being able in spite
of everything to preserve the old. Of both opposites it holds that they have
learned nothing and forgotten nothing. When he who is filled with an ideal
also possesses intellectual honesty, he will, for the ideal's own sake,
investigate all the conditions under which it is to work, and all the ways
in which this working must intervene in human affairs.

The matter now stands thus, that the fundamental values and the norms
themselves come into being only under quite definite historical conditions.
Ideals which should have arisen and been formed without the influence of
historical realities have never existed. Plato's ethical idealism did not
form itself in his consciousness without the influence of the past of Greek
intellectual life, and his ethics has, in spite of all opposition to his
people's custom and to his contemporaries' notions, a decidedly Greek stamp.
Christianity has in its original form a decidedly Oriental stamp, and it
points back to Jewish tradition and to Persian notions. Kant's ethics has a
long and significant history of development. If I am here also to speak of
myself, I have expressly emphasized that the presuppositions on which I have
built my presentation of ethics are due to historical development and cannot
be said to have been present at all times. When
\textsc{Westermarck}\footnote{\textit{Origin and Development of Moral
Ideas}. I, p.~18.} says that if the word ``ethics'' is to be used as the
name of a science, its task can only be to study the moral consciousness as
a fact, I do not feel myself hit by this remark, which is among other things
also directed against me; for I have always assumed a
historical-psychological development of the moral consciousness and
maintained that every ethics will correspond to a certain form and a certain
stage of this process of development. (See Etik III, 13; 21). This does not
deprive ethical ideas of their significance or their validity. One can after
all steer by the stars, although they have a history. And besides, the
ethical voyage takes place on earthly seas and is therefore conditioned not
merely by the stars, but by wind and wave.

The task must always be a development of what is historically given in the
direction of the ideal norms. But there must at the same time take place an
orientation with regard to whether these norms are indeed rightly set up,
and whether they are indeed really the expression of experiences of value.
If for one reason or another the fundamental value itself is shaken, it
becomes a psychological-historical question whether transition to another
fundamental value is possible. Here too there is the possibility of serious
crises and tragic catastrophes.

\subsection{Ethics and the Individual Relations.}
\label{sec:4d}
% pp. 90--92
% VARIANT READING recorded in the transcription: the printed body head reads
% "Etiken og de individuelle Relationer", the Indhold on p. 109 "Etik og de
% individuelle Relationer" (without the -en). Both were verified on the
% images; the printed body head is what is followed here.

\noindent Besides the psychological relations, which concern the
acknowledgment of a fundamental value and of the norms following from it,
and the historical relations, which concern the possibility of carrying
these norms through under given conditions, there comes forward in a deeper
ethical discussion yet a third kind of relations, which concern the
individual's possibility of satisfying the norms. If ethical demands are not
to hover in the air, they must be set in such a way that the individual's
capacity and need are taken into consideration. Only then will a striving in
the direction of what is demanded be able to be released. There cannot
therefore be mechanically demanded the same of everyone. Of him who is
capable of great or self-sacrificing activity there is demanded something
other and more than of him who under inner discord or by reason of
powerlessness of mind perhaps just manages to keep his head above water.
Only then is the same ethical work in reality demanded of all; for ethical
work is the expression of the energy with which inner and outer resistance
is overcome, and if this resistance is different for different individuals,
the work actually performed can be determined only by the relation to this
different resistance. What is demanded of the individual may then lie now
above, now below the average demands set by a general law. Everyone is, as
Schiller's Wilhelm Tell says, taxed according to ability.

In a treatise on the laws of relation in ethics (Etiske Undersøgelser, 1891)
I have in detail grounded that individualization of ethical norms which is
as necessary as it is a difficult consequence of the concept of relation. It
rests on the principle that the law exists for man's sake, not man for the
law's sake. The work the individual must do in order to satisfy ethical
norms, the virtues and the duties demanded of him, must themselves serve as
means for his personal development, and this presupposes that his capacity
and need become co-determining for the demands which according to the norm
are to be made of him in particular.

Most ethical theories have misjudged the significance of this individual
relation. Nevertheless it has again and again asserted itself. In theology
it lies at the basis of the concept grace, in jurisprudence of the concept
pardon. For these concepts are unethical when regard for the individual
conditions does not lie at the basis of them. In philosophy the analogous
concept equity has been set up, which stands in opposition to abstract
justice, but in reality, if it is an ethical concept at all, must be said to
express the highest justice, since it takes the individual moments into
consideration. In the most recent time's individualization of the execution
of punishment this consideration also comes forward; what a transgressor of
the law is sentenced to is precisely the definite execution of punishment
actually applied to him.

From the significance of the individual relations there follows an important
consequence for social ethics. It must be a necessary task to strive to
provide all members of society with such possibilities of development that
the choice of work can become free. So long as this is not attained, the
ethical demand upon society has not been satisfied. On this point the social
struggle will more and more come to be joined.

\subsection{Kant's Ethics and the Concept of Relation.}
\label{sec:4e}
% pp. 92--94

\noindent Kant, who himself, as we have seen above, saw in the concept of
relation a fundamental law for all thinking, would nonetheless not
acknowledge it in ethics. The ethical law shall according to him in no
respect or degree have proceeded from experience; psychological and
historical relation shall thus not hold for it. And it shall demand the same
of all, whatever individual presuppositions they may have. But for that
reason the ethical law stands, according to Kant, also as an inexplicable
marvel before the working, thinking and striving human race, which is
constantly placed before definite conditions. Kant states this quite openly.
``How pure reason can be practical, human reason is quite incapable of
explaining, and all trouble, all labour to seek an explanation of it is
wasted''. (Grundlegung\textsuperscript{3}, p.~125). ``How a law can in and
for itself and immediately be a determining ground for the will, that is for
human reason an insoluble problem''. (Kritik der praktischen Vernunft,
p.~128). Kant does indeed concede that the law can work in us only by
calling forth a feeling of the law's own sublimity; but he then declares
this feeling to be psychologically inexplicable.

But from this relationlessness in psychological, historical and individual
respects it again follows that we can at bottom pass no properly ethical
judgments. When we have no access to understanding the relation between
norm, motive and action, we can never have the right to call any human
action before us good or evil. This consequence Kant already drew in the
first edition of ``Kritik der reinen Vernunft'' (p.~551), though only in a
note below the text: ``Die eigentliche Moralität der Handlungen bleibt uns
daher, selbst die unseres eigenen Verhaltens, gänzlich
verborgen''.\footnote{Perhaps the statement of \textsc{Droysen} going in the
same direction (see above p.~81) is in part due to the influence of Kant's
train of thought.} This is a correct consequence when the concept of
relation is not to be applicable in ethics. The matter stands otherwise when
ethical norms are conceived as the expression of endeavours which are
precisely formed into more perfect forms through the necessary relations
asserting themselves in all thinking and acting. Precisely the concept of
relation opens great possibilities for us here. It demands that we set the
one compass-leg at a definite place in the real world (in a psychological
fact, in something historically given, in an individual possibility) and
then try how far we can carry the other compass-leg out. If we cannot at
once set the other compass-leg where we wish, we can perhaps gradually move
it further out. Perhaps even the first compass-leg can be moved. The point
in ethics as in physics is definite points de repère.

Just as space and time are ``Ideas'' (in Kant's sense of the word), and just
as experience (reality) is an ``Idea'', so too the ethical good is an
``Idea'', which acquires real significance not by excluding all relations,
but by using them to set the special tasks and to judge the attempts at
their solution.

Kant and Droysen are nonetheless not entirely wrong in their declarations
that the proper morality of actions cannot be discovered or grounded. For we
can never be sure that we have taken all relations into consideration. It is
at last impossible to satisfy the demand which precisely in virtue of the
concept of relation must be made, namely to win insight into the relation
between the action's inner and outer conditions and into the magnitude of
the ethical work which the individual man has performed or can perform from
his individual conditions. All ethical judgments are therefore imperfect and
hold at most only approximately.

Kant's assumption of a norm holding for all men can also still be
maintained. For the principle of relation becomes at bottom one with the
principle of personality, that is, the principle that every man shall always
be regarded and treated as end, never merely as means. This is satisfied
when the demands upon the individual are not merely socially motivated, but
also individualized according to the individual's capacity and need, so that
his personality can be developed through the work he does to follow the
norm.

Two of Kant's most significant ethical thoughts thus come into their full
right only through application of the concept of relation.

\subsection{Philosophy of Religion and the Concept of Relation.}
\label{sec:4f}
% pp. 94--98

\noindent The religious problem too, however one may stand towards it,
exhibits a many-sided application of the concept of relation. This problem
rests in all its forms on an opposition between the valuable on the one
side, the valueless or value-hindering on the other, --- more briefly
expressed, between value and reality, in that reality exhibits both value
and its opposite. The opposition here present need not --- in my view --- be
a dualism, as has been thought\footnote{\textsc{La Harpe}: \textit{La
Religion comme conservation de la valeur, dans ses rapports avec la
philosophie générale de Harald Høffding} (1920), p.~93 ff.}; but it may
certainly rise to a catastrophe, bring with it that tragedy of life the
possibility of which no view of life can exclude, as little as one can
exclude the possibility of earthquakes. And even if in the tragic, in the
suffering that leads to the limit of the human, one would find the true
nature of existence expressed, that would precisely confirm the significance
of oppositions for the conception of life. The tragic presupposes relation
in the sharpest sense, without however excluding a holding fast to the
valuable. The tragic hero maintains humanity's nobility in the greatest
catastrophe. But apart from the great catastrophes too, the relation between
value and reality appears as the problem on which religion under all forms
rests. In all forms of religion the point is what is victorious in the
world, the good or the evil. Within its great fullness reality exhibits the
possibility of development in various directions, and the contention between
forces working in the direction of what men ascribe highest value to, and
other forces working in a quite other, perhaps opposite direction, sets
precisely the questions which religion in one form or another seeks to
answer. It is this contention that leads the religious consciousness to form
two concluding concepts, God and world, and all religious thinking is a
brooding over the relation between these two concepts. In asceticism and in
mysticism one seeks, now more in practice, now more in speculation, to reach
beyond all relations, or (as we have seen above p.~47 Note) at least to make
the one term of the relation disappear; but again and again relation forces
itself forward and sets the problem anew. The need of life disturbs the
ascetic, and the world of reality forces its way into the mystic's cell with
its manifoldness and its unrest. One discovers at last the position of the
one compass-leg, however well it was hidden.

The question is whether the valuable in the world (the valuable part of
reality) is a foreign guest which has no home there, or whether it should
not rather be a disclosure of existence's inmost nature. The only rational
answer which can here be given (and which I have attempted to give in
``Totalitet som Kategori'' Kap.~5) must proceed from the close connexion
between the concepts value and whole. Everything to which value is ascribed
is the expression or the condition of a certain whole's subsistence and
development, and the valuable arises according to the same laws that make
formations of wholes (for example organisms, personalities, societies)
possible. (Cf.\ above p.~83 f.). When now whole stands against whole, or
when chaotic elements make the formation of wholes impossible, then there
are present precisely those experiences which lead to a distinction between
value and reality.

When the monotheistic religions regard God as the origin of all reality, the
problem of evil appears with particular sharpness. Outside the origin of all
reality, evil cannot have its origin. Therefore bold and profound
speculation, such as Jacob Böhme's, has drawn the consequence that there
must be a dark ground in God himself, out of which what conflicts with value
in the world has arisen. Inner relations in God are then set in place of the
outer relations at which experience stops. Still more complicated does the
matter become when, as Böhme already recognized, it may be a condition for
the development of the good that evil should exercise its resistance and
thereby work towards the release of forces which would otherwise lie fallow.
A consideration of this kind led Goethe (Aus meinem Leben. Fjerde Del) to
emphasize what he called the demonic in the world, and to use the motto:
nemo contra deum nisi deus ipse.\footnote{See further on this
\textit{Religionsfilosofi}, §~88---89.}

Religious faith is a hope or an assurance that there must always in
existence (in that which experience exhibits, or in another) be the
possibility of valuable formations of wholes. It is this that is expressed
by saying that religion is faith in the persistence of value. But just as
wholes arise and perish, so too do values. The religious consciousness
therefore comes in the course of its development to recognize that it cannot
be the sensuously or empirically valuable that persists, indeed that it is
at last impossible to give a concept or an image of what it is that
persists. This recognition does not work its way forward without strong
resistance and demands a deep resignation. It is easily veiled through the
imagery religion employs.

The concept God itself contains a valuation; men's God is what they set
highest. In the expression ``divine'' this is plainly seen; it is a
determination of value. Plato already speaks of ``that which makes God
divine'', which is that God lives constantly in the great Ideas, which men
can do only now and then (Phaedrus p.~249 C). Plotinus goes a step further,
in saying that what makes the gods divine must itself be the highest
divinity (Ennead. V, 1, 2). The continuation of this train of thought would
lead all theology back to a doctrine of value.\footnote{Cf.\ with this the
way in which the expression ``divine'' is used by a modern mystic
(\textit{Oplevelse og Tydning}, p.~89).} ---

The expression ``the persistence of value'' says nothing about how much
value there is in the world, or of what kind it is. On these two points lie
the differences between the various religions or religious standpoints,
which this is not the place to enter upon; that must be left to the history
of religion and the psychology of religion. For the theory of relation it
was here only the task to show that the concept of relation involuntarily
asserts itself in all religious thinking, the most elementary as well as the
most speculative. It is, as we saw, partly the relation between value and
reality, partly the relation between eternal and perishable values, that
sets the questions for religious thinking. New relations arise during the
historical development of the religions and often push the underlying
relations into the background. Thus the relation between a religion's proper
kernel and its formation in cult and dogma. In its zeal to obtain guarantees
for the persistence of value, the religious consciousness has a tendency to
multiply the object of faith, so that the one dogma is used to support the
other. --- A particularly important relation is that between religion and
morality. Consciousness may be so taken up with the persistence of value and
with the dogmas thought to guarantee it, that energy is drawn away from the
work of discovering and bringing forth values or of maintaining the values
already found. One may dwell so long at the will's boundary that one does
not come to work in the domains where men need not wait for the gods. And at
last, from an ethical point of view, the consideration forces itself forward
that value is nonetheless not necessarily dependent on persistence. One can
live in and for values even if no certainty can be won about their
persistence. The divine is not dependent on time. What is valued as beautiful
and good can keep its value, whatever its fate may become.


% ============================================================
%  5. THE CONCEPT OF RELATION IN COSMOLOGY
% ============================================================

\section{The Concept of Relation in Cosmology.}
\label{sec:5}
% pp. 99--108 (body ends two thirds down p. 108; p. 109 carries the Indhold)
%
% Section 5 is divided by BARE LATIN LETTERS a. b. c. d. (pp. 99, 103, 105,
% 107), set in ordinary roman at the paragraph indent with the text running
% straight on -- no title, no run-in dash, no letterspacing, no line of their
% own. They are lettered paragraphs, not heads, and the Indhold lists no
% sub-entries under section 5 at all. So, as in the transcription, they are
% kept inline with no sectioning command and no \addcontentsline: there is no
% title anywhere in the book to give them.

Does the carrying through of the concept of relation permit a rounded, even
if not concluded, conception of existence? Will not every conception of a
whole here show itself impossible in face of a chaos of points of view, each
of which lays claim to justification? --- In face of such questions there
are various considerations to bring forward.

a.~No orientation is possible except from a definite standpoint. But every
standpoint can in turn be drawn into the discussion, because it stands in
definite relations and rests on definite presuppositions. There is a
constant tendency to overlook this. Sound common sense, the religious
world-views, and all too often natural-scientific and philosophical systems
as well, have proceeded on the assumption that there are standpoints whose
grounding cannot in turn be demanded, and whose justification cannot be
discussed; what is the basis for a world-view is itself not to be
conditioned. Every world-view has its primal phenomenon, to use Goethe's
expression, and Goethe expressly asserted that the primal phenomenon cannot
be explained, but must be regarded as an absolute last. But wherever one may
flee, the concept of relation will be able to assert itself, and there may
arise that unrest which is inseparable from all working thought. The idylls
of thinking are always exposed to dissolution and can only be provisional
resting places. In the most recent time one is beginning to become clear
that a rational explanation of all parts or sides of existence from a single
primal phenomenon is not possible. But one then often maintains that here at
last a choice must be made between two possibilities: between a
mathematical-mechanical and a psychological-qualitative primal phenomenon.
Natural science approaches more and more the mathematical ideal of a world
of identities or substitutions, where differences of time and quality play
no part, while the human sciences (or, as some prefer to call them, the
sciences of culture) maintain the reality of differences of quality and
stages of development which exclude reduction to identity. Between the two
types thus determined the choice is supposed to lie. In this direction Wundt,
James Ward and Henri Bergson have expressed themselves. But that such a
choice should be necessary would presuppose that an either---or is here
present, that there thus exist no other types of existence than the two
indicated. The circumstance that identity, quality and temporality cannot be
reduced to one does not justify an either---or. Irreducibility is one thing,
absence of all other possibilities another. In face of an assertion of an
either---or there is established here the significance of Spinoza's
remarkable doctrine of the infinitely many attributes under which existence
appears, but of which our experience shows us only two --- precisely those
two types between which the choice is supposed to lie. To a far higher
degree than he is conscious of, Spinoza here points to the limitation of our
knowledge.\footnote{\textit{Spinoza's ``Ethica''}, p.~22---24.} When
Spinoza, in spite of the manifoldness of points of view whose possibility
and justification he acknowledged, nonetheless maintained a doctrine of
unity, it was because he too had his primal phenomenon, which he was
convinced could not be denied from any of the many points of view. He held
to the fact that we have a science --- that rational understanding of the
world is possible to the extent that it is. It was this fact that he
projected outward as an absolute reality and called the substance in
everything. And it is at any rate a foothold making it possible to discuss
the various points of view that may assert themselves, particularly also
with regard to whether their mutual relation is contradictory or merely
contrary. The very limits to which human knowledge is always subject can be
determined only by scientific thinking. And here it shows itself, as
remarked several times in the foregoing, that the necessity of relations can
be both a barrier and a spur for knowledge. Where we can find no definite
relations, we have no knowledge either. Our seeking therefore does not
cease; we know what we lack: new positions for our compass-legs. Therefore
we end with a questioning.

It has now even been thought that the carrying through of the concept of
relation would make all truth impossible. Bergson distinguishes between
relative and limited knowledge; the first distorts the nature of its object,
the second leaves it untouched, in that it grasps only a part of
it.\footnote{\textit{Bulletin de la socièté française de Philosophie.} VIII,
p.~341.}
% The Danish note prints "socièté" with a grave on the first e; logged as a
% printer's defect at its site in transcription.tex and kept as printed.
But with the relations the problems too fall away; there become neither
definite questions to put nor definite answers to give. Consistently
Bergson therefore points to an immediate intuition as the only way by which
truth can be attained and the inmost essence of existence be given. In the
closer carrying through of this, however, he makes use of constant
inferences from analogy --- and to that extent of relations too, since
analogy means likeness of relation.\footnote{\textit{Henri Bergson's
Filosofi}, p.~54.} An example of how relation forces its way in precisely
when one has thought to be able to exclude it! --- The most important
objection to Bergson's philosophy is from a philosophical point of view not
the one Poincaré has somewhere expressed in the sentence: Le monde
Bergsonien n'a pas des lois! --- but this, that this philosophy has no
problems. It is really a glorification of the problemless state.

The life of human thought expresses itself in the relations it is able to
bring out and to determine. Every such determination means a victory over a
chaos. In a chaotic series of difference no definite relations whatever are
demonstrable; absolute interchangeability prevails. That the bringing out
and determining of relations is thought's essential property follows, as
mentioned several times in the foregoing, from relation being correlative to
synthesis. The unity and energy of thought, of consciousness, express
themselves in the relations that let the individual objects appear in their
peculiarity. A comprehending together by which no relations are set is a
contradiction. Even in dreams the law of determination through relation
prevails, however much they may approach chaos.

Nor is it in and for itself an imperfection in our knowledge that at
different times and in different places, according to the different points
of departure, it comes to different results. Thereby there is merely set a
new problem, a new task for knowledge, namely to find out how the different
points of departure and thereby the different results became possible. In
the foregoing we have seen how the physicist and the historian recognize the
possibility of different standpoints and of the different conceptions
following from them. It is a task which of course cannot always be solved,
but which necessarily presents itself, since truth must be able to explain
both itself and the errors; that belongs to complete truth. Plato already
saw this when he (Republic X, p.~602 C---603 A) shows that when different
observers apprehend the same thing in different ways as regards magnitude
and figure, one can by counting and measuring arrive at a conception from
which the various individual conceptions can be explained. Galileo was
inspired by this thought in the founding of his doctrine of motion. But he
was embarrassed by what his imperfect telescope showed him of the phenomenon
now called Saturn's ring. The phenomenon appeared now as edges or corners on
the planet, now as two different subsidiary planets. At last Galileo gave up
all inquiry about it. Later Huyghens with his improved telescopes saw
plainly that Saturn was surrounded by a ring, and showed how earlier
conceptions could be explained from this.\footnote{\textsc{Hoefer}:
\textit{Histoire de l'Astronomie}, p.~445---447.} It is the triumph of
thought in every domain when a ``perfect'' knowledge can be won from which
the ``imperfect'' conceptions can be explained. When Spinoza teaches (Eth.
II, 36. Dem.) that our ``imperfect'' presentations arise with the same
necessity as the ``perfect'' ones, it is precisely by the ``perfect'' (i.e.
more perfect) presentation that this necessity is recognized. That is the
meaning of Spinoza's famous proposition: veritas est norma sui et falsi
(Eth. II, 43. Schol.).

Thought thus sets its boundary problems in virtue of its own law, and it is
in virtue of this same law that it explains its errors.

b.~But the very counting and measuring to which Plato pointed, and of which
Galileo made use in principle, is after all due to human mental activity,
and what right have we to ascribe validity to its results? We wish after all
to understand the world as it is, and not as it presents itself to human
calculation, even where the latter attempts to adopt ever so essential
points of view.

Already above (p.~58 ff.) I have, in face of such a consideration,
impressed that the concept of relation must be applied not merely between
objects among themselves, but also between the objects and the knowing
subject, or, more correctly expressed, that the knowing subject is itself an
object. ``The percipient event'' too, to use Whitehead's expression, belongs
among the events in the world. We have no right to assert in advance an
absolute opposition between object and subject. Even if one seeks to
transform all science into formal science and reckons everything that cannot
be brought into this form among ``anthropomorphisms'', these anthropomorphisms
are after all a fact which also demands its scientific investigation.

At the basis of those questions there really lies the popular and dogmatic
definition of truth as our knowledge's agreement with reality. This static
concept of truth is, as has often been shown, self-contradictory. Reality
does not stand fully finished, so that we need only go and view it. It first
comes into being for us through strict mental labour in observation and
thinking. The history of this labour is the history of truth. It serves to
characterize existence itself that for a certain stage of the development of
human knowledge it had to appear in a certain definite way. We need only,
for example, enter into all the presuppositions of the Middle Ages'
rejection of antipodes in order to see what truth had to stand as necessary
at this stage. What is reality for men at a given stage depends on how far
the mental labour at this stage has advanced. Later this reality may stand
as an imperfect first attempt. The eternity of truth rests on the continuity
of mental labour.

It is not an intellectually aristocratic train of thought to which we are
hereby led. Precisely according to the concept of relation there is truth in
every honest striving after clarity. We do not ourselves create the
presuppositions from which we think, and we do not even become fully
conscious of them. There is here an analogy between the intellectual and the
ethical life (see above p.~81 ff.). In the necessity with which the
``imperfect'' presentations appear, eternal truth is itself present. And
often it has even been a condition for progress in the knowledge of truth
that full earnest was made of the (as it later showed itself) false
presuppositions. As William Hamilton has said: a living untruth is better
than a dead truth.

Beyond the constant striving to which the dynamic concept of truth points we
do not get. Our thought cannot leap over its own shadow. Our last foothold
is the working life of thought itself in its constant discovering and
determining of relations. A conclusion is here not possible.

Bacon and Comte maintained that man is to be explained by the world, not the
world by man. But here there is no dilemma. We have, as we have seen (above
p.~60 ff.), never an absolute subject over against an absolute object, but
indeed an objectively determined subject (S\textsubscript{o}) over against a
subjectively determined object (O\textsubscript{s}), and they both demand a
searching characterization. And when we then compare them, this comparison
must be undertaken by a subject, so that the problem constantly recurs. He
who would explain man from the world must not forget that the world is known
only through human inquiry, and that all investigation of the relation
between world and man is instituted by a human thought.

c.~When we believe that existence is more comprehensive than what our
thought at any time or in any form can comprehend, this is grounded in our
experiencing again and again how new relations and thereby new questions
emerge and demand determinations. We have therefore, as shown above, no
right to believe that existence is exhausted by the two primal phenomena or
attributes which can briefly be denoted as mind and matter. It is the
teasing thing about the concept of relation that it constantly shows us a
hole in our systems. Even if we have formed the most beautiful totalities in
our thought, the question emerges what this totality stands in relation to.

Investigation of the concept of totality in its significance for the
world-view ends, as I have elsewhere (Totalitet som Kategori, Kap. 6) sought
to show, in two problems: that of the possibility of saying anything about
existence as a whole, and that of the possibility of determining the
position of the valuable within such a whole --- thus in a cosmological and
in a religious problem. When the demand for grounding is made sharply and
definitely, these problems can be solved only by the way of analogy. Of this
a critical investigation of cosmological and religious notions also
testifies. (Cf.\ the epistemological part of my Religionsfilosofi.) For the
philosopher such analogies will lie nearest as are drawn from the working
life of thought and will, whose history is for us the history of truth and
goodness. It is due to a mystical need in Spinoza when he projected outward
the connexion of thought which for him was the primal phenomenon as a
resting substance. Our last analogies must rely on our most significant
experiences, and these have labour, striving and struggle for their content,
corresponding to the constant contention between our two most fundamental
categories, synthesis and relation. Without the appearance of oppositions
and striving to overcome them, existence cannot, philosophically seen, be
thought. Among these oppositions the opposition between the valuable and the
value-hindering asserts itself, practically seen, most strongly.

If no grounded assurance can be won that all oppositions and inhibitions can
be harmonized, the analogy with the working life of thought and will can
nonetheless afford a poetic and symbolic expression for the assurance that
the labour in which we are placed is not an isolated feature in reality; for
that, relation and the struggle with relations hang far too closely together
with the knowledge of existence which we are at all able to win.

d.~Only in the resting form of the image can we, for a while, get beyond the
constant unrest which by reason of the concept of relation is characteristic
of every waking thought. It is therefore also to the world of images that he
who seeks rest and peace takes his refuge. A famous Danish theologian has
stated that without the use of images we should come out into the void, in
that we should have to hold to some poor concepts that could give no light.
Yes, the philosopher has only his poor concepts to hold to. But if these
concepts under constantly renewed reflection show themselves to be grounded
in the nature of the life of thought and of experience, then he may have the
right to believe that the pulse of existence beats in the best and clearest
thoughts men can form, and perhaps does not beat most faintly where new
questions are set in virtue of the concept of relation. And he reserves his
judgment on the analogies lying at the basis of the formation of images in
the individual case. As little as any thought will any image be able
completely to express the inconcludable existence of which we are members.
In the fact that every parable limps, the concept of relation announces
itself again and again.

In Obstfelder's ``En Præsts Dagbog'' the struggle between thought and image
is depicted in a profound and atmospheric way, as it is often waged in a
yearning mind. The poet is clear --- as the philosopher too can be --- that
the highest dwells ``further within us than we ourselves can penetrate''.
But he is not satisfied with this conviction. He wishes to have it out
before him, to ``come outside that which is within us'', to be able to say
thou to it, to have an image of it. If he had succeeded in reaching this
goal, he would have made the experience that no image suffices, but that a
constant need for new images asserts itself, just as the investigator
constantly discovers new relations by which his hitherto formed concepts are
burst asunder.

The world of poetry stands open to all, also to the philosopher. But the
philosopher cannot acknowledge any image formed once for all. And he will be
clear that when he has passed over into the world of poetry, the thoughts
fall silent, just as for the Anacreontic poet ``the cares fall silent'' when
he grasps his cup. There is constantly required new labour, both with images
and with thoughts, in order that life may come into its own. Perhaps what an
Egyptian sage called ``a yearning mind, whose words are all hidden'' is the
most inward and highest form of the life of the spirit.

% ============================================================
%  THE PRINTED INDHOLD (p. 109), translated.
%  Reproduced as a page, following the transcription's editorial choice to
%  preserve what the book itself prints rather than generate it. (The
%  \tableofcontents at the front of this file is a translator's convenience
%  and is a different thing.) Layout as printed: heading, short centred rule,
%  the column head "Side" at the right, twenty-four entries with dot leaders,
%  a second rule, and the two dating lines.
%
%  Letterspacing: only the five NUMBERED entries and the word INDHOLD are
%  spærret in the printing; the lettered and Greek sub-entries are ordinary
%  roman. The five are set \emph{} here, as in the transcription.
%
%  Divergences preserved, neither reading normalised:
%    * the Greek sub-entries of 3b and 3c are titles that exist ONLY here --
%      the printed pages carry bare marks with no title text.
%    * entry 4d reads "Etik", where the printed head at p. 90 reads "Etiken".
%    * section 5's four lettered paragraphs are not listed at all.
%    * the Indhold's entries carry no terminal full stop; the printed heads do.
% ============================================================
\clearpage
{\centering\large CONTENTS\par}
\vspace{0.6ex}
{\centering\rule{2.4cm}{0.4pt}\par}
\vspace{1.2ex}
\begingroup
\parindent=0pt
\setlength{\parskip}{0pt}
\small
\mbox{}\hfill Page\par
\hangindent=1.6em \emph{1.~The History of the Concept of Relation}\dotfill 3\par
\hangindent=1.6em \emph{2.~The Concept of Relation in Psychology}\dotfill 16\par
\hangindent=1.6em \emph{3.~The Concept of Relation in Epistemology}\dotfill 26\par
\hangindent=3.2em \hspace*{1.6em}a.~The Concept of Relation and the Other Fundamental Categories\dotfill 26\par
\hangindent=3.2em \hspace*{1.6em}b.~The Concept of Relation and the Formal Categories\dotfill 33\par
\hangindent=4.8em \hspace*{3.2em}α.~Pure Quality-Series\dotfill 33\par
\hangindent=4.8em \hspace*{3.2em}β.~Identity and Relation\dotfill 40\par
\hangindent=4.8em \hspace*{3.2em}γ.~The Relations of Negation and Rationality\dotfill 44\par
\hangindent=3.2em \hspace*{1.6em}c.~The Concept of Relation and the Real Categories\dotfill 45\par
\hangindent=4.8em \hspace*{3.2em}α.~Causality, Reality and Relation\dotfill 45\par
\hangindent=4.8em \hspace*{3.2em}β.~The Relations of Totality and Development\dotfill 55\par
\hangindent=3.2em \hspace*{1.6em}d.~The Relation Subject---Object\dotfill 58\par
\hangindent=4.8em \hspace*{3.2em}α.~Introduction\dotfill 58\par
\hangindent=4.8em \hspace*{3.2em}β.~On Copernicanism\dotfill 62\par
\hangindent=4.8em \hspace*{3.2em}γ.~Philosophical Standpoints\dotfill 67\par
\hangindent=4.8em \hspace*{3.2em}δ.~The New Theory of Relativity\dotfill 72\par
\hangindent=1.6em \emph{4.~The Concept of Relation in Ethics}\dotfill 77\par
\hangindent=3.2em \hspace*{1.6em}a.~Historical Research and the Concept of Relation\dotfill 77\par
\hangindent=3.2em \hspace*{1.6em}b.~Ethics and the Psychological Relations\dotfill 81\par
\hangindent=3.2em \hspace*{1.6em}c.~Ethics and the Historical Relations\dotfill 87\par
\hangindent=3.2em \hspace*{1.6em}d.~Ethics and the Individual Relations\dotfill 90\par
\hangindent=3.2em \hspace*{1.6em}e.~Kant's Ethics and the Concept of Relation\dotfill 92\par
\hangindent=3.2em \hspace*{1.6em}f.~Philosophy of Religion and the Concept of Relation\dotfill 94\par
\hangindent=1.6em \emph{5.~The Concept of Relation in Cosmology}\dotfill 99\par
\endgroup
\vspace{5ex}
{\centering\rule{2.4cm}{0.4pt}\par}
\vspace{2ex}
{\centering\small Presented at the Meeting of 21 October 1921.\par
Completed at the Printing House 14 November 1921.\par}

\end{document}
